% Generated by scripts/make_ms_version.py from paper_tmlr/main.tex. Do not edit by hand:
% edit the TMLR source and rerun the script.
\documentclass[MS,12pt]{ucdavisthesis}
../paper_tmlr/math_commands.tex
\usepackage{natbib}
\usepackage{hyperref}
\hypersetup{hidelinks}
\usepackage{url}
\usepackage{graphicx}
\usepackage{booktabs}
\usepackage{multirow}
\usepackage{amsmath}
\usepackage{xcolor}
\usepackage{placeins}
\usepackage{flafter}
\usepackage{float}
\usepackage{longtable}
\usepackage{adjustbox}
\usepackage[htt]{hyphenat}
\makeatletter\setlength{\@fptop}{0pt}\makeatother
\renewcommand{\topfraction}{0.9}\renewcommand{\bottomfraction}{0.8}\renewcommand{\textfraction}{0.07}
\renewcommand{\floatpagefraction}{0.75}\setcounter{topnumber}{3}\setcounter{bottomnumber}{2}\setcounter{totalnumber}{5}
\renewcommand{\bibfont}{\singlespacing}
\setlength{\bibsep}{4pt plus 2pt}

\newcommand{\task}[1]{\texttt{#1}}
\newcommand{\pending}[1]{\textcolor{red}{[#1]}}

\title{Re-Tokenization, Segmentation, and a Blind Spot: What Content-Free Boundary Markers Do in Character-Level Transformations}
\author{Faezeh Khazaee}
\authordegrees{}
\officialmajor{Computer Science}
\graduateprogram{Computer Science}
\degreemonth{December}
\degreeyear{2026}
% Major professor (chair) and committee members are left blank on purpose.
\committee{}{}{}{}{}
% Plan II project report: the document-type line reads Project (change to Thesis for Plan I).
\makeatletter\def\@Thesisname{Project}\makeatother

\abstract{Inserting a content-free marker (\texttt{|}) into the input every $k$ characters of a character-level transformation, and asking the model to copy it through, is a cheap single-call intervention. We study what it does across seven tasks, seven open instruction-tuned models, and two versions of Gemini Flash (3.5 and 3.6), using two controls that fix tokenization at one token per character (verified at the open models, supported by API token counts at Gemini 3.6): the same letter tasks with spaced letters, and exact relabelings with digits. We find three separate effects. First, on letter strings, markers act largely as re-tokenization: at five of the seven open models, spacing the letters with no markers improves accuracy more ($+0.16$ to $+0.33$) than markers do ($+0.13$ to $+0.26$). Second, the gain survives one token per character: with spaced letters or digits, markers remove 33\% to 83\% of the remaining errors at the four strongest open models, persist to $N = 256$ at Qwen2.5-32B, and match there with commas or line breaks as markers. Third, with a marker every two or more characters, the position right after a marker fails whenever its output depends on the input before it: on four dependency tasks it falls to near chance at the stronger models while the next position recovers, and rewriting the tasks so that position needs nothing before the marker repairs it. Blocking marker-attending heads removes part of the gain, and a constrained-decoding ablation on four tasks at Qwen2.5-32B shows that marking either input or output alone hurts while marking both helps; the cross-boundary failure arises from input markers alone. Real promoter DNA and English text show the same gain and failure. In practice: separate characters before adding markers, mark long sequences sparsely on input and output, and keep markers away from positions whose output depends on what came just before them.}

\begin{document}
\makeintropages

\chapter{Introduction}
\label{sec:intro}

Ask a language model to apply a simple, fully specified rule to every character of a string (shift each letter forward, substitute through a lookup table, output the XOR of each digit with the one before) and accuracy does not fade gracefully with length. It holds near ceiling on short inputs and then collapses well before any obvious capacity limit \citep{hou2025focus}. Standard remedies change the computation itself: splitting the input across several calls \citep{zhou2023leasttomost}, asking for intermediate steps \citep{nye2021scratchpad, chen2023pot}, or adding memory to the architecture \citep{dai2019transformerxl, bulatov2022rmt, mohtashami2023landmark}. Each costs extra calls, extra tokens, or training.

We study a cheaper intervention. A content-free marker character (\texttt{|}) is inserted into the input every $k$ characters, and the model is told to copy it through to the output unchanged. The marker carries no task information, needs no extra call, and requires no training. On letter tasks it raises accuracy substantially, by up to $+0.26$ on average at the stronger models we test. Our question is why, and when it hurts.

A marker does more than mark a boundary. Tokenizers merge runs of letters into multi-letter tokens, and a marker after every few letters splits those runs, so on letter strings a marker is also a tokenization intervention. We separate the two with two controls that keep one token per character whether or not markers are present: the same letter tasks with a space between letters, and exact relabelings of them with digits. We find three effects (Figure~\ref{fig:overview}):

\begin{enumerate}
\item \textbf{On letter strings, markers largely re-tokenize} (Chapter~\ref{sec:tokenization}). At all five Qwen and Gemma models, spacing the letters with no markers at all helps more than markers on unspaced letters ($+0.16$ to $+0.33$ against $+0.13$ to $+0.26$); at Mistral and Phi the two help about equally. The striking accuracy of the position right after a marker on letter tasks (positive in all 27 task--model cells, with per-model means of $+0.07$ to $+0.34$ over the no-marker run at the same positions) is largely a token-boundary effect: with spaced letters or digits that position is no better than the rest of its segment.
\item \textbf{The gain persists when each character is its own token} (Chapter~\ref{sec:segmentation}). With spaced letters or digits, markers remove 33\% to 83\% of the remaining errors at Qwen2.5-14B and 32B, Qwen3.8-27B, and Gemma-4-31B (significant in 23 of 24 pooled cells with generation-level standard errors; at $N = 64$ all 24 cells are positive, $+0.07$ to $+0.37$, and 21 of 24 clear $|z| \ge 2$ under a cell-level bootstrap; Appendix~\ref{app:stats}), and the two controls agree closely. The gain is absent at short lengths, where these models are near ceiling, appears at $N = 32$, and at $N = 64$ reaches $+0.27$ (spaced letters) and $+0.22$ (digits), with little dependence on spacing. At Qwen2.5-32B it persists to $N = 256$ (56\% of remaining errors removed) and is the same with commas or line breaks as markers. At the three smallest models the gains are smaller and inconsistent (significant in 7 of 18 cells).
\item \textbf{The position after a marker fails on dependencies that cross it} (Chapter~\ref{sec:blindspot}). With a marker every two or more characters, on tasks where each output also depends on the previous input, the position right after a marker falls below the no-marker run in 33 of 35 task--model cells (every cell at the six strongest models), to near chance at the stronger models, even where the model is otherwise near ceiling. The next position recovers immediately, and rewriting the tasks so that this position needs nothing before the marker repairs it. The loss is not created by how the instruction describes the markers: telling the model they are only visual separators and that the rule continues across them narrows it, and removes it only at one model on two of the four tasks, and commas or line breaks in place of \texttt{|} leave the position at chance up to $N = 256$. Pooled task-level effects on these tasks are consistent with a mix of this loss and the within-segment gains of Finding~2, which is why their sign varies with length and model.
\end{enumerate}

We also test the obvious mechanism (Chapter~\ref{sec:mechanism}). In Qwen2.5, some heads attend from each position to the adjacent marker, but at Qwen2.5-32B and $N = 64$ blocking the ten most marker-attending heads removes 1\% to 19\% of the gain (9\% rather than 13\% on spaced letters once a single outlier trial is set aside), against at most 5\% for controls, and blocking 200 removes 21\% to 28\%, against none for controls of the same size: marker-directed attention contributes, but it is not concentrated in the heads that attend to markers most. A $2 \times 2$ separating input from output markers under constrained decoding, on four tasks at Qwen2.5-32B (Section~\ref{sec:sides}), finds that marking either side alone \emph{hurts} ($-0.03$ to $-0.16$) while marking both helps ($+0.02$ to $+0.39$); the blind spot follows the input markers alone. Chapter~\ref{sec:practical} turns the findings into guidance.

\paragraph{Supporting analyses.} Table~\ref{tab:summary} lists the checks behind each finding, from alignment-robust scoring and other delimiters to base checkpoints, activation patching, and a finetuned separator, and Chapter~\ref{sec:real} repeats the design on real DNA and English text; none changes the sign of any of the three effects.

\begin{figure}[!htb]
\centering
\includegraphics[width=\linewidth]{figures/fig0_overview.pdf}
\caption[The three effects at a glance.]{The three effects at a glance. (a)~Qwen2.5's tokenizer merges runs of letters into multi-letter tokens; a marker every four letters splits only some of them, while spacing the letters gives one token per letter. (b)~With one token per character (spaced letters or digits), markers still help the four strongest open models, from $N = 32$ on (the marker runs, ``Bounds'', minus the no-marker runs, ``Base''; mean over models and the three position-local tasks; whiskers are 95\% intervals; same data as Figure~\ref{fig:profile}a). (c)~On \task{running\_xor}, each output is the XOR of its input and the one before; the position right after a marker, whose dependency crosses the marker, falls to near chance at the stronger models, and the next position recovers.}
\label{fig:overview}
\end{figure}

% =====================================================================
\FloatBarrier
\chapter{Related Work}
\label{sec:related}

\paragraph{Tokenization, formatting, and position.} Subword tokenizers obscure the characters inside a token: models know only partially which characters their tokens contain \citep{kaushal2022tokens}, character-level manipulation is a known weakness \citep{edman2024cute}, and longer tokens obscure where a character sits \citep{uzan2026charbench}. Models handle non-canonical encodings of the same string surprisingly well, splitting text into characters can improve string manipulation \citep{zheng2025broken}, retokenizing a prompt while preserving its bytes probes segmentation sensitivity \citep{jain2026retokenization}, and instruction-tuned models reconstruct canonical tokens internally from character-split input \citep{yang2026wordrecovery}, which is one reason our spaced-letter control works at all. Surface formatting can change few-shot accuracy by tens of points \citep{sclar2024quantifying}, token costs vary across languages and formats \citep{petrov2023unfairness, ahia2023cost}, and the surface form of numbers \citep{nogueira2021arithmetic}, positional description \citep{shen2023positional}, and the presentation of positions \citep{duan2023interpolation} change arithmetic accuracy and length generalization. These works change how text is encoded. We hold the encoding fixed instead: our controls give one token per character whether or not markers are present, so a marker's remaining effect cannot be a change in how characters are grouped into tokens, and the post-marker failure can be attributed to the boundary rather than to the tokenizer.

\paragraph{Learned segmentation.} Tokenizer-free models learn where to place boundaries: H-Net learns content-dependent chunks over bytes end to end \citep{hwang2025hnet}, H-Net++ adds a cross-chunk attention module because chunking alone ignored dependencies across chunks \citep{zakershahrak2025hnetpp}, and MrT5 learns to delete byte tokens partway through the encoder \citep{kallini2025mrt5}. Byte-level models still segment the byte stream: ByT5 reads raw bytes \citep{xue2022byt5}, MegaByte groups them into fixed patches \citep{yu2023megabyte}, and the Byte Latent Transformer into entropy-based patches \citep{pagnoni2024blt}. Our markers are the opposite case, fixed content-free boundaries inserted at inference time into a model never trained to use them. Whether the first position after a learned or fixed patch boundary also loses access to what came before it is open; the H-Net++ design suggests that dependencies across boundaries need explicit support there as well.

\paragraph{Attention sinks and separator tokens.} Language models concentrate attention on a few tokens, most famously the first, which act as sinks \citep{xiao2024streaming, gu2024sink, zhang2025catch}; \citet{su2026sinksurvey} survey uses and interpretations. SepLLM shows that separator tokens can summarize the segment before them \citep{chen2024sepllm}, SinkLoRA exploits sinks during fine-tuning \citep{zhang2024sinklora}, punctuation tokens have served as anchors for trained sparse attention \citep{qiu2026phsa}, and tabular structure in a prompt improves factual data interactions \citep{oh2024tables}. Closest to our intervention, SepSeq \citep{sun2026sepseq} inserts a separator every $k$ items into long numerical inputs at inference time, without training, to counter the dispersion of softmax attention, and improves sequence-level answers such as counts and maxima; there only the input is segmented and the answer is a single value. Our tasks are per-position transformations whose output is written segment by segment, and there we separate re-tokenization from segmentation, find that segmenting the input alone hurts unless the output is segmented too (Section~\ref{sec:sides}), and find that the position after a separator loses the dependency across it. Those separators are learned or exploited in training. Ours are inserted at inference time into models never trained for this use, so any effect must come from the boundary rather than from task content, although the bar may carry structural associations from pretraining, which is one reason we also test commas and line breaks. \citet{dangal2026anchoring} find that inserting a comma at a clause boundary reduces degradation of later predictions while biasing attention toward existing punctuation does not, because the inserted token brings a new hidden state computed from the full context; our marker is such a token, and the weak effect of blocking attention to it fits that reading.

\paragraph{Attention heads.} Induction heads implement in-context copying \citep{olsson2022induction}, and causal head gating assigns roles to heads by intervention \citep{nam2025headgating}. We characterize the marker-attending heads and block them from the markers. On character-counting probes, character-level information is encoded in early layers and then actively suppressed by late-layer components, with no attention heads specializing in routing it \citep{datta2026suppression}; our finding that the marker benefit is not concentrated in the heads that attend to markers most is consistent with that picture.

\paragraph{Long sequences, decomposition, and benchmarks.} Models use the middle of long contexts poorly \citep{liu2024lost}, make sparse but systematic errors on state tracking \citep{liu2023flipflop}, fail long-horizon tasks for reasons that compound \citep{wang2026mirage, pushkin2026lead, anand2026periodic}, and show an accuracy cliff on repetitive deterministic tasks of the kind we use \citep{hou2025focus}. Remedies range from architecture, such as SWAN's interleaved global and sliding-window layers \citep{puvvada2025swan}, ALiBi \citep{press2022alibi} and FIRE \citep{li2024fire}, Transformer-XL \citep{dai2019transformerxl}, recurrent memory \citep{bulatov2022rmt}, landmark tokens \citep{mohtashami2023landmark}, and learned context management \citep{kontonis2026memento}, through inference-time reshaping of attention or rotary frequencies \citep{han2024lminfinite, wang2024resonance}, to prompting that decomposes the computation \citep{zhou2023leasttomost, chen2023pot, nye2021scratchpad}; content-free markers sit at the cheap end of this spectrum, and the tokenization-independent benefit we find is a long-sequence effect of this kind. Boundaries matter for probabilistic conditioning as well: byte-level sampling methods exist because a prompt that ends mid-token distorts the next-token distribution \citep{hayase2025bytesampler}, and the failure we document is a deterministic, per-position counterpart at an inserted boundary. Our tasks are generated fresh for every trial, following BeyondBench \citep{srivastava2025beyondbench}; external checks use ProcBench \citep{fujisawa2024procbench}, the copying task of \citet{jelassi2024repeat}, and CLRS-Text \citep{markeeva2024clrstext}.

% =====================================================================
\FloatBarrier
\chapter{Experimental Setup}
\label{sec:setup}

This section describes the tasks and the three formats they are posed in, the marker intervention, the models and their decoding settings, and the metrics and statistics used throughout.

\section{Tasks and Formats}

Every task maps an input string of length $N$ to an output string of the same length under a deterministic rule stated fully in the prompt (Table~\ref{tab:tasks}). In \emph{position-local} tasks, output $y_i$ depends only on input $x_i$. In \emph{one-step dependency} tasks, $y_i$ also depends on $x_{i-1}$; in the \emph{windowed} task, on $x_{i-2}, x_{i-1}, x_i$. For dependency tasks, the first position is copied as-is. Inputs are sampled uniformly with a deterministic seed per (task, length, spacing, alphabet size, trial), so every model sees the same inputs; because the seed includes the spacing, different spacings, including the no-marker run, draw independent inputs from the same distribution, so main-grid comparisons are unpaired (a paired-input rerun agrees to within 0.01; Appendix Table~\ref{tab:paired}). \task{copy\_previous} requires no computation, only a one-step lookback, so any harm from markers cannot be blamed on arithmetic; \task{sliding\_window\_max} uses a three-symbol window, so its output stays informative at every length.

The position-local tasks are posed in three formats. In the \emph{letter} format (the default), letters are written contiguously, and tokenizers merge them: at each of the seven open models, whose tokenizers we can inspect, 48 random letters encode to 23 to 32 tokens without markers and to one token per letter with a marker after each letter. In the \emph{spaced} format, the same letters are separated by single spaces and the model is asked to answer in the same format; the 48 letters encode to exactly 48 tokens at every open model whether or not markers are present, and each marker adds one token of its own. In the \emph{digit} format, the letters are relabeled as digits ($A \to 0$, $B \to 1$, \dots), which every open-model tokenizer encodes one per character. On Gemini 3.6 Flash, whose tokenizer we cannot inspect, the API's token counts show exactly one token per added character and one per marker in both controls, against 0.46 tokens per contiguous letter (Appendix Table~\ref{tab:geminitokens}). Both controls are exact transformations of the letter tasks: same seeds, same positions, same rule and prompt wording (verified in every cell). The dependency tasks use digits throughout. One caveat on \task{conditional\_shift}: its rule depends on a symbol's position in the alphabet, which for digits equals the digit itself, so its digit version is easier.

\begin{table}[H]
\centering
\small
\caption[Tasks. $A$ = alphabet size, $M$ = modulus.]{Tasks. $A$ = alphabet size, $M$ = modulus. Lengths $N \in \{4, 8, 16, 32, 64\}$. Position-local tasks are run with letters, spaced letters, and digits; dependency tasks use digits.}
\label{tab:tasks}
\adjustbox{max width=\linewidth,max totalheight=0.86\textheight}{\begin{tabular}{llll}
\toprule
Task & Tier & Rule & Axis \\
\midrule
\task{cyclic\_shift} & position-local & next symbol in the alphabet, wrapping & $A \in \{2,4,8\}$ \\
\task{random\_permutation} & position-local & fixed arbitrary substitution (a derangement) & $A \in \{2,4,8\}$ \\
\task{conditional\_shift} & position-local & symbols at even alphabet positions forward, odd back & $A \in \{2,4,8\}$ \\
\midrule
\task{copy\_previous} & one-step & $y_i = x_{i-1}$ & binary \\
\task{running\_xor} & one-step & $y_i = x_i \oplus x_{i-1}$ & binary \\
\task{running\_sum\_local} & one-step & $y_i = (x_i + x_{i-1}) \bmod M$ & $M \in \{2,4,8\}$ \\
\task{sliding\_window\_max} & windowed & $y_i = \max(x_{i-2}, x_{i-1}, x_i)$ & $A \in \{4,8\}$ \\
\bottomrule
\end{tabular}}
\end{table}

\section{The Marker Intervention}

In the marker condition (``Bounds''), a \texttt{|} is inserted after every $k$ input characters, with $k \in \{1, 2, 4, 8, 16\}$; only $k < N$ is used. The expected output is the ground truth computed on the continuous, marker-free sequence, with markers spliced in at the same positions. The prompt states the rule, then instructs, for example: \emph{``The input contains vertical bars (|). These are structural boundaries. DO NOT transform them. Copy them exactly to the output.''} (the verb matches the task, e.g.\ ``shift'' for \task{cyclic\_shift}). The no-marker baseline (``Base'') uses the identical prompt, including that instruction, so that the only difference between conditions is whether markers appear in the input. That instruction is odd for an input with no bars in it, so we measured what it costs: rerunning the no-marker condition with every mention of markers removed, on the same inputs, across 300 (task, length, alphabet) cells at the four stronger models, raises accuracy by $+0.006$ on average (95\% interval $[0.002, 0.009]$), by at most $+0.026$ in any model--family combination, and the sentence helps rather than hurts in 105 of the 300 cells (Appendix~\ref{app:stats}). At Qwen2.5-7B, rechecked for the review, it costs $-0.003 \pm 0.009$. The marker effects we report are therefore overstated by well under a percentage point on average against a baseline whose prompt never mentions bars. Example prompts are in Appendix~\ref{app:prompts}. Runs also included $k = 3$, which we exclude from the pooled power-of-two aggregates but analyze separately, with $k = 5$ and 6, as a check on periodicity (Chapter~\ref{sec:blindspot}, Appendix Table~\ref{tab:blindk}).

\section{Models and Decoding}

Open models are Qwen2.5-7B/14B/32B-Instruct \citep{qwen2024qwen25}, Mistral-7B-Instruct-v0.3 \citep{jiang2023mistral}, Phi-3.5-mini-instruct \citep{abdin2024phi3}, Qwen3.8-27B \citep{qwen2026qwen38}, and Gemma-4-31B-it \citep{gemma2026gemma4}, run locally in bfloat16 with each model's chat template and no system message beyond the template's own default (thinking disabled for Qwen3.8 and Gemma 4) and greedy decoding. Runs use PyTorch 2.6.0 (CUDA 12.4) and Transformers 5.12.1 with the eager attention implementation, on NVIDIA A100 80GB GPUs (Titan RTX for some smaller-model runs); Appendix~\ref{app:trials} lists the exact checkpoint revisions. Qwen3.8-27B is a hybrid architecture: 48 of its 64 layers use Gated DeltaNet linear attention \citep{yang2025gated} and only 16 use softmax attention; Gemma-4-31B-it uses softmax attention in every layer (50 sliding-window and 10 global layers). The generation budget is $1.5N + 20$ new tokens, or $2N + 20$ at $k = 1$, where a correct marked answer is $2N - 1$ tokens; predictions are clipped to the answer's length before scoring, and among the 86,400 stored raw generations a budget cut-off affects at most 0.6\% of any condition class and 0.04\% at the four stronger models, nearly all of it Phi-3.5-mini on spaced letters (Appendix~\ref{app:trials}, which also lists which batches store raw output).\footnote{The $k = 1$ cells at $N \ge 64$ were first generated under the smaller budget, which cut them off, and then regenerated; with greedy decoding the new outputs extend the old ones, except for 480 of the 5,820 regenerated outputs, nearly all of Mistral, Phi, and Qwen2.5-7B on \task{conditional\_shift} and \task{random\_permutation} (and a few of Phi on \task{running\_xor} and Qwen2.5-7B on \task{copy\_previous}), whose original runs used a different GPU model; these diverge after a shared prefix of about 20 characters (median) with the same accuracy over it (0.375 against 0.373), and dropping them changes no pooled effect by more than 0.002.} The closed models are two versions of Gemini Flash, queried through the API with the same prompts and seeds and with thinking disabled: Gemini 3.5 Flash \citep{gemini35flash} at temperature 0 on the letter and dependency tasks with 3 trials per cell, and, after Google began routing its requests to a newer model \citep{geminichangelog}, Gemini 3.6 Flash \citep{gemini36flash} on all 13 task--format variants with 5 trials per cell. Gemini 3.6 ignores the temperature parameter, so its outputs are samples rather than greedy decodes. Appendix~\ref{app:trials} gives the query dates and the checks on the served model and its usage metadata. Qwen2.5-32B is our primary model for head-level experiments.

\section{Metrics and Statistics}

\emph{Character accuracy} is the fraction of output positions whose character matches the ground truth, after stripping markers and spaces from both sides and clipping the prediction to the ground-truth length. We avoid the name ``token accuracy'' deliberately: the unit is a character of the answer, not a tokenizer token. Marker placement is not part of this metric; Appendix~\ref{app:wfr} reports exact transformation match and marker-placement compliance separately.

\paragraph{Pooled effects.} Only spacings $k < N$ are run, so marker runs are concentrated at longer lengths than no-marker runs, and accuracy falls with length. A pooled effect therefore compares the mean character accuracy of the marker runs with the no-marker mean reweighted to the same length distribution, $\Delta = \overline{\mathrm{acc}}_{\text{Bounds}} - \sum_N p_N\, \overline{\mathrm{acc}}_{\text{Base},N}$, where $p_N$ is the share of marker runs at length $N$; it pools $N \le 64$, $k \in \{1, 2, 4, 8, 16\}$, and alphabet or modulus values, and $z$ is $\Delta$ over its standard error. Because the controls leave strong models near ceiling, we also report the \emph{share of remaining errors removed}, $\bar\Delta / (1 - \overline{\mathrm{acc}}_{\text{Base}})$: the mean $\Delta$ over a model's tasks divided by one minus the mean length-matched no-marker accuracy. Pooled over $N \le 64$, it includes the short lengths where $\Delta \approx 0$, so it is conservative relative to the share removed at $N = 64$, which Chapter~\ref{sec:segmentation} reports separately. Open models use 20 to 80 trials per cell (Appendix~\ref{app:trials}); Gemini 3.5 uses 3 and Gemini 3.6 uses 5. We call a pooled cell significant at $|z| \ge 3.54$, the Bonferroni-corrected two-sided $p < 0.05$ across all 125 pooled cells we test (the 77 in the overview of Chapter~\ref{sec:blindspot} and 48 for the two controls).

\paragraph{Units of uncertainty.} A marker run and the no-marker run it is compared with see different random inputs from the same distribution, so comparisons are unpaired; a paired rerun at four models shows that pairing would have cut standard errors by 9\% to 22\% and moved no estimate by more than 0.01 (Appendix Table~\ref{tab:paired}). Standard errors are $s/\sqrt{n}$ over generations, combined as $\mathrm{SE}(\Delta) = \sqrt{\mathrm{SE}_{\text{Bounds}}^2 + \sum_N p_N^2\,\mathrm{SE}_{\text{Base},N}^2}$. A bootstrap over generations changes the Bonferroni verdict in 1 of 91 cells; one over whole experimental configurations leaves 16 of 91 cells above the threshold instead of 62, so we read pooled significance as conditional on the grid we ran and rely on the length- and position-resolved analyses for the structure of the effects (Appendix~\ref{app:stats}).

\paragraph{Redundant configurations.} At $A = 2$ the three position-local rules coincide: the only derangement of two symbols, the cyclic successor, and the conditional shift all flip the bit. At $M = 2$, \task{running\_sum\_local} computes \task{running\_xor}. These cells still test different \emph{instructions} for the same function, but they are not independent evidence across tasks; Appendix~\ref{app:stats} recomputes every pooled effect with them dropped, which moves the median $|\Delta|$ by 0.013, makes the effect larger in 54 of 91 cells, and leaves 23 of the 24 Finding~2 cells significant.

\paragraph{Per-position accuracy.} For a marker run with spacing $k$, output position $i$ (0-indexed, markers stripped) is \emph{post-marker} if $i > 0$ and $i \bmod k = 0$, and \emph{within-segment} otherwise. Each class is compared with the no-marker run's accuracy at the same position indices, which controls for accuracy varying along the sequence. Per-position results pool $N \le 64$, $k \in \{2, 4, 8, 16\}$, and alphabet or modulus values; $k = 1$ has no within-segment positions, so per-position claims concern spacings of two or more. The unit for uncertainty is the generation, and intervals are 95\%.

% =====================================================================
\FloatBarrier
\chapter{Tokenization Effects on Letter Strings}
\label{sec:tokenization}

\begin{figure}[!b]
\centering
\includegraphics[width=\linewidth]{figures/fig_three_formats.pdf}
\caption[Character accuracy without markers (left end of each line) and with markers (right end, with its 95\% interval) on the three position-local tasks, one line per model and one panel per format.]{Character accuracy without markers (left end of each line) and with markers (right end, with its 95\% interval) on the three position-local tasks, one line per model and one panel per format: contiguous letters, spaced letters, and digits. Mean over tasks; $N \le 64$, $k \in \{1, 2, 4, 8, 16\}$; no-marker values are length-matched to the marker runs (Chapter~\ref{sec:setup}). At the Qwen and Gemma models, spacing the letters without markers beats markers on contiguous letters. With tokenization fixed, markers still help the stronger models.}
\label{fig:formats}
\end{figure}

Figure~\ref{fig:formats} compares the three formats. On contiguous letters, markers raise accuracy at every model: by $+0.26$ at Qwen3.8-27B, $+0.24$ at Qwen2.5-32B, $+0.23$ at Gemma-4-31B, $+0.20$ at Qwen2.5-14B, $+0.13$ at Qwen2.5-7B, $+0.07$ at Mistral-7B, $+0.06$ at Phi-3.5-mini, and $+0.22$ at both Gemini versions (all 27 cells positive, 25 significant). But at the five Qwen and Gemma models, and at Gemini 3.6, simply spacing the letters, with no markers, helps more: $+0.33$ at Qwen3.8, $+0.32$ at Gemma 4 and Gemini 3.6, $+0.30$ at Qwen2.5-32B, $+0.27$ at Qwen2.5-14B, and $+0.16$ at Qwen2.5-7B. At these models, spaced letters without markers are more accurate than contiguous letters with markers (for example, 0.80 against 0.74 at Qwen2.5-32B, 0.89 against 0.79 at Gemma 4, and 0.97 against 0.87 at Gemini 3.6); at Mistral ($+0.06$ from spacing) and Phi ($+0.05$), spacing and markers help about equally. Without markers, the letter tasks are hard for these models largely because of how letters are tokenized: averaged over lengths, the four stronger models solve the spaced-letter versions at 0.70 to 0.89 and the digit versions at 0.78 to 0.93, against 0.44 to 0.56 with contiguous letters (0.81 to 0.94 on digits for the two tasks whose digit version is an exact relabeling, leaving out \task{conditional\_shift}, whose digit rule is easier). Qwen2.5-7B is the exception, weaker on digits than on letters, because it writes the letters t or o in place of 0 in 34\% to 39\% of its digit outputs (Appendix~\ref{app:models}).

\begin{figure}[!b]
\centering
\includegraphics[width=\linewidth]{figures/fig1_per_position.pdf}
\caption[Per-position effects, every model.]{Per-position effects, every model. For each model, the change in accuracy with markers at the within-segment positions (orange circles) and at the position right after a marker (blue diamonds), each against the no-marker run at the same position indices, with 95\% intervals; one panel per block, mean over its tasks; $N \le 64$, $k \in \{2,4,8,16\}$. Left: on contiguous letters the post-marker position gains the most. Middle: with spaced letters it gains no more than the rest of its segment. Right: on the dependency tasks it loses at every model while the within-segment positions hold or gain (Chapter~\ref{sec:blindspot}). Gemini 3.5 was not run on spaced letters.}
\label{fig:position}
\end{figure}

The per-position view (Figure~\ref{fig:position}) shows how. On contiguous letters, the position right after a marker beats the no-marker run at the same positions in all 27 task--model cells, with per-model means from $+0.07$ (Phi) to $+0.34$ (Qwen2.5-32B, Qwen3.8, and Gemini 3.6), and it beats the within-segment positions of the same marker runs in 25 of 27 cells. A marker guarantees a token boundary right before the following letter: at $k = 1$ that letter becomes its own token, and at $k \ge 2$ it starts a new token, which may still span several letters (Figure~\ref{fig:overview}a). With spaced letters or digits, where every symbol is already its own token, that advantage is gone: in raw marked accuracy the post-marker position beats the within-segment positions in only 12 of the 24 task--model cells of each control (the seven open models and Gemini 3.6), and the two classes are within 0.04 of each other at every model, with the post-marker position slightly the lower of the two at Qwen2.5-7B; measured as the change against the no-marker run at the same positions, as in Figure~\ref{fig:position}, they are within 0.07 (Appendix Table~\ref{tab:perpos}). That is pooled over spacings; at the sparsest spacing, $k = 16$, the post-marker position on the controls falls 0.07 to 0.16 below the position after it at Qwen2.5-7B, Mistral, and Phi, and on spaced letters at Qwen2.5-14B, a small echo of the dependency-task loss that the stronger models do not show. The spacing profile tells the same story (Figure~\ref{fig:profile}b, next section): on contiguous letters at $N = 64$, the gain is largest with a marker after every letter ($+0.47$ at the four stronger models) and falls steadily as markers get sparser ($+0.09$ at $k = 16$), as expected if markers help by splitting letter runs into shorter tokens. And the letter gain is present even on four-letter strings ($+0.15$), where the controls show no marker effect at all.

% =====================================================================
\chapter{Tokenization-Independent Gains on Long Sequences}
\label{sec:segmentation}

\begin{figure}[!t]
\centering
\includegraphics[width=\linewidth]{figures/fig_profile.pdf}
\caption[The marker effect (Bounds $-$ Base) at the four stronger models (Qwen2.5-14B and 32B, Qwen3.8-27B, Gemma-4-31B), mean over the three position-local tasks.]{The marker effect (Bounds $-$ Base) at the four stronger models (Qwen2.5-14B and 32B, Qwen3.8-27B, Gemma-4-31B), mean over the three position-local tasks. (a) By sequence length, pooled over spacing. (b) By marker spacing at $N = 64$. Whiskers are 95\% intervals over the 12 model--task pairs; the three tasks show the same pattern. The two tokenization controls (spaced letters, digits) agree: no effect up to $N = 16$, a large effect at $N = 64$, and little dependence on spacing. Contiguous letters show a different signature: a gain at every length that is largest for the densest markers.}
\label{fig:profile}
\end{figure}

\begin{figure}[!t]
\centering
\includegraphics[width=\linewidth]{figures/fig_length_spacing_grid.pdf}
\caption[Character accuracy over the full length-by-spacing grid, one panel per model.]{Character accuracy over the full length-by-spacing grid, one panel per model. Rows: marker spacing (``none'' is the no-marker run); columns: sequence length; color: accuracy, mean over the tasks of the block and their alphabet or modulus values; blank cells have $k \ge N$. Top: the three position-local digit tasks, where the marker rows stay bright as $N$ grows while the no-marker row fades at the stronger models. Bottom: the four dependency tasks, where the rows with $k \ge 2$ are darker than the no-marker row at short lengths and for dense spacings, while sparse spacings recover or gain at $N = 32$ and 64 at the stronger models.}
\label{fig:grid}
\end{figure}

Holding each character to one token does not hold the whole prompt fixed: markers still add tokens, lengthen the sequence, and shift every later position. What the controls establish is narrower, that the gain is not an artifact of merging characters into multi-character tokens; we call it \emph{tokenization-independent} in that sense throughout. With tokenization fixed, markers still help the stronger models. With spaced letters, the gain is $+0.13$ at Qwen2.5-32B, $+0.11$ at Qwen2.5-14B, $+0.09$ at Qwen3.8, and $+0.06$ at Gemma 4, removing 62\%, 38\%, 55\%, and 48\% of their remaining errors; with digits it is $+0.06$ to $+0.09$, removing 33\% (Qwen2.5-14B) to 83\% (Gemma 4), where the share of remaining errors removed is $\bar\Delta / (1 - \overline{\mathrm{acc}}_{\text{Base}})$ over the model's three tasks. Because these shares pool $N \le 64$, including the short lengths where $\Delta \approx 0$, they understate the share removed at $N = 64$. Leaving out \task{conditional\_shift}, whose digit rule is easier than its letter rule, raises rather than lowers them (to 44\% to 98\% with spaced letters and 55\% to 96\% with digits), because it is the hardest of the three tasks in both formats. These effects are significant in all 12 spaced-letter cells and 11 of 12 digit cells at these four models (generation-level standard errors; the one exception is \task{conditional\_shift} with digits at Qwen2.5-14B), and a stratified estimator that treats the length, spacing, and alphabet cells as fixed strata gives the same 23 of 24. At $N = 64$ alone, all 24 cells are positive ($+0.07$ to $+0.37$) and all clear $|z| \ge 3.54$ at the generation level, while the more conservative cell-level bootstrap leaves 21 of 24 above $|z| \ge 2$ and 4 above 3.54 (Appendix Table~\ref{tab:controlcells} lists every cell). At the three smallest models the gains are smaller and inconsistent: with spaced letters, $+0.02$ (Qwen2.5-7B), $+0.06$ (Mistral), and $+0.03$ (Phi); with digits, $+0.05$, $-0.01$, and $+0.10$; significant in 7 of these 18 cells. Gemini 3.6 Flash is at 0.94 to 1.00 on both controls without markers, so they leave it no room to show the effect: markers add $0.00$ to $+0.04$ there, none significant; in the one control cell where it is below 0.95 without markers, spaced letters at $N = 64$ (0.90), they add $+0.09$.

This benefit is a long-sequence effect (Figure~\ref{fig:profile}a; Figure~\ref{fig:grid} shows the whole length-by-spacing grid at every model). Up to $N = 16$, the four stronger models are near ceiling in both controls (0.92 to 0.96 without markers) and markers change nothing. At $N = 32$, accuracy without markers starts to fall and markers add $+0.07$ to $+0.09$; at $N = 64$, where accuracy without markers is 0.59 (spaced letters) and 0.68 (digits), markers add $+0.27$ and $+0.22$. In the strongest case, Qwen2.5-32B on \task{random\_permutation} with eight digits at $N = 64$, accuracy rises from 0.41 to between 0.87 and 0.99 with any spacing from every 2 to every 16 characters. The gain depends little on spacing (Figure~\ref{fig:profile}b): at $N = 64$ it is $+0.21$ to $+0.32$ for spaced letters and $+0.17$ to $+0.26$ for digits across all spacings, slightly weaker at the extremes (a marker after every character, which nearly doubles the number of output tokens, or after every 16).

The gain is spread over the whole segment rather than concentrated after the marker (Figure~\ref{fig:position}, middle). A similar tokenization-independent gain appears inside segments on the dependency tasks, which use digits: within-segment positions gain $+0.07$ at Qwen2.5-32B, $+0.05$ at Gemma 4 and Gemini 3.5, and $+0.02$ at Qwen3.8 and Gemini 3.6 over the no-marker run at the same positions, though not at Qwen2.5-14B. Under edit-distance alignment these within-segment gains shrink but stay positive at the stronger models (for example, at Qwen2.5-32B the gain drops from $+0.11$ to $+0.03$ on \task{running\_xor} and from $+0.11$ to $+0.05$ on \task{running\_sum\_local}, and at Qwen3.8 from $+0.13$ to $+0.06$ and from $+0.14$ to $+0.09$), so part of what markers do inside segments is keep the model from slipping out of alignment. The gap between the two scorings is the alignment-sensitive part of the gain, the part that index-by-index scoring attributes to fewer alignment slips; it is 52\% of the within-segment gain in the median cell among the 11 of 16 dependency-task cells at the four stronger models whose within-segment gain is positive (range 14\% to 85\%; Appendix Table~\ref{tab:aligned}). The remainder survives alignment and so reflects fewer substitutions, though the split between two scoring rules is a descriptive decomposition, not a causal one. Markers also change how errors cluster along the sequence (Appendix Table~\ref{tab:autocorr}): at $N = 64$ under index scoring, the probability that an error is followed by another error is 0.70 to 0.77 without markers on the spaced and digit formats at the four stronger models and 0.24 to 0.59 with a marker every 2 to 16 characters, and the mean run of consecutive errors shortens from 3.1 to 3.9 positions to 1.3 to 2.3; under alignment the same probability moves only from 0.34 to 0.51 to 0.21 to 0.40, so what sparse markers mostly shorten is the cascade after a slip. A marker after every character lowers the error rate but leaves that probability at 0.45 to 0.70. On the dependency tasks the same table shows a slight rise in the chance of a fresh error with sparse markers; that rise is the blind spot's isolated errors (22\% to 55\% of errors there are single errors between correct positions), and excluding pairs that touch a post-marker position brings it to or below its no-marker level, while the drop in the chance that an error follows an error survives the exclusion. The gain also extends beyond the lengths of our main grid: at Qwen2.5-32B, on \task{random\_permutation} with digits ($A = 8$, $k \in \{4, 16\}$), markers remove 89\% to 91\% of the remaining errors at $N = 64$, 78\% to 81\% at $N = 128$, and 56\% to 57\% at $N = 256$ (accuracy 0.18 without markers, 0.64 to 0.65 with), the ranges spanning \texttt{|}, commas, and line breaks as markers (Appendix Table~\ref{tab:delimiter}). One reading is that a marker every few characters helps a strong model keep its place in a long sequence; Chapter~\ref{sec:mechanism} states that hypothesis precisely and tests the simplest version of it. The experiments in this section establish the effect, not its cause. Chapter~\ref{sec:mechanism} shows that the heads that attend most to markers carry little of this effect, and Qwen3.8-27B, whose layers are mostly linear attention, shows it comparably to the other strong models.

% =====================================================================
\chapter{The Post-Marker Blind Spot on Dependency Tasks}
\label{sec:blindspot}

Figure~\ref{fig:poscurve} shows the pattern along the sequence: at Qwen2.5-32B, on a position-local task every position gains, while on \task{running\_xor} the accuracy curve drops at exactly the positions that follow a marker and recovers at the next, and the same columns light up at every model. Figure~\ref{fig:examples} shows what this looks like in single generations: with a marker every four characters, the errors of an otherwise correct output sit at the position after each marker. On the four dependency tasks, with a marker every two or more characters, the position right after a marker is below the no-marker run at the same positions in 33 of 35 task--model cells (Figure~\ref{fig:position}, bottom; these tasks use digits, so tokenization plays no role). The two exceptions are slight ($+0.02$ each, Qwen2.5-7B on \task{running\_xor} and Phi on \task{sliding\_window\_max}); every cell at Qwen2.5-14B and 32B, Qwen3.8, Gemma 4, and both Gemini versions shows the loss. We call this the \emph{blind spot} for short. The term describes the behavior, not a claim about what the model represents: accuracy near chance shows that the information before the marker is not used correctly here, not that it is inaccessible to the model. The loss grows with model strength: $-0.03$ to $-0.06$ at Phi, Mistral, and Qwen2.5-7B; $-0.15$ at Qwen2.5-14B and 32B; $-0.19$ at Qwen3.8; $-0.26$ at Gemma 4; and $-0.43$ and $-0.34$ at Gemini 3.5 and 3.6, which are near ceiling without markers. At the stronger models the post-marker position is close to chance: on \task{running\_xor} (chance 0.5) it is correct 0.51 to 0.55 of the time at Qwen2.5-14B and 32B, Qwen3.8, Gemma 4, and both Gemini versions. The failure is confined to the positions whose inputs lie across the marker. On the one-step tasks that is one position: on \task{running\_xor} at $k \ge 4$ the change against the no-marker run is $-0.15$ at the position right after a marker and $+0.20$ at the next, averaged over the four stronger models. \task{sliding\_window\_max} needs the two previous inputs, so the first \emph{two} positions after a marker cross it: pooled over $k \in \{2, 4, 8, 16\}$ at the four stronger models, the first changes by $-0.13$, the second by only $+0.03$, and the strictly internal positions gain $+0.08$ (Appendix Table~\ref{tab:blindwindow}). We count that second position as within-segment throughout, which dilutes that task's within-segment gain and leaves the second crossing step out of the post-marker estimate; at $k = 2$ every within-segment position of \task{sliding\_window\_max} is such a second position. With a marker every 16 characters, accuracy one position later is back at 0.85 (Qwen2.5-32B), 0.70 (Qwen3.8), 0.82 (Gemma 4), 1.00 (Gemini 3.5), and 0.91 (Gemini 3.6). A marker after every character ($k = 1$) is different: there every position follows a marker, and on \task{copy\_previous} at $N = 32$ Qwen2.5-32B is more accurate with such markers (0.88) than without (0.83), plausibly because a fully regular pattern makes the lookback fixed (Chapter~\ref{sec:mechanism}).

\begin{figure}[!t]
\centering
\includegraphics[width=\linewidth]{figures/fig_position_curve.pdf}
\caption[Accuracy by output position ($N = 32$).]{Accuracy by output position ($N = 32$). (a, b) Qwen2.5-32B, mean over trials and alphabet or modulus values with 95\% bands: on a position-local digit task every position gains from markers and the no-marker run decays along the sequence; on \task{running\_xor} the position right after each marker (dots) drops to about chance while the positions between markers gain, and with markers every 8 characters the drops move to positions 9, 17, and 25. (c, d) Every model with a marker every 4 characters: the change in accuracy at each position relative to the no-marker run at the same position, pooled over the three position-local digit tasks (c) and the four dependency tasks (d). The dependency panel shows the blind spot as the orange columns at positions 5, 9, 13, and so on at every model; Gemini 3.5 was run on fewer tasks.}
\label{fig:poscurve}
\end{figure}

\begin{figure}[!t]
\centering
\includegraphics[width=\linewidth]{figures/fig_examples.pdf}
\caption[Three generations of Qwen2.5-32B at $N = 32$, character by character.]{Three generations of Qwen2.5-32B at $N = 32$, character by character: the input as shown to the model, the expected output, and the model's output, with correct characters in green and wrong ones in red; bars are the markers. Top: \task{running\_xor} without markers. Middle: the same task with a marker every 4 characters, where the errors sit at the position right after each marker (arrows) and nowhere else. Bottom: a position-local task (\task{random\_permutation} with eight digits) with the same markers, where every position is right.}
\label{fig:examples}
\end{figure}

\paragraph{Pooled effects on dependency tasks.} Averaged over all positions (Figure~\ref{fig:split}), a dependency task's response combines this post-marker loss, on about a fraction $1/k$ of positions, with the within-segment gain of Chapter~\ref{sec:segmentation} on the rest, so its sign varies with length and model: of 35 dependency cells, 8 are significantly positive and 7 significantly negative, five of them on \task{copy\_previous}, which is negative at seven of nine models (Appendix~\ref{app:perpos}).

\begin{figure}[t]
\centering
\includegraphics[width=\linewidth]{figures/fig1_split_heatmap.pdf}
\caption[Pooled change in character accuracy when markers are added (Bounds $-$ Base), contiguous-letter position-local tasks and digit dependency tasks, every model.]{Pooled change in character accuracy when markers are added (Bounds $-$ Base), contiguous-letter position-local tasks and digit dependency tasks, every model. Pooled over $N \in \{4,\dots,64\}$, $k \in \{1,2,4,8,16\}$, and alphabet or modulus values, with the no-marker mean length-matched to the marker runs. $^*$: $|z| \ge 3.54$ (Bonferroni across the 125 pooled cells we test). Blue: markers help; red: markers hurt; ---: not run. Bottom, set apart: dependency tasks rebuilt so the position after each marker needs nothing before it (Section~\ref{sec:reset}); these rows change the rule at the post-marker position and their prompts shift the no-marker baseline, so they are not comparisons on the original tasks, and Appendix Figure~\ref{fig:reset} gives the per-position test.}
\label{fig:split}
\end{figure}

\section{Robustness of the Blind Spot}
\label{sec:robust}

The checks in this subsection change how the effect is measured or delivered and leave the blind spot in place; Table~\ref{tab:summary} collects them.

\paragraph{Alignment-robust scoring.} Character accuracy compares characters index by index, so a dropped or inserted character misaligns everything after it. Re-scoring every run with an edit-distance alignment, which repairs such slips instead of counting the rest of the output as wrong, leaves the post-marker loss in place at every strong open model and task (Appendix Table~\ref{tab:aligned}). At Qwen2.5-32B it is $-0.24$ on \task{running\_xor} and $-0.22$ on \task{running\_sum\_local} under alignment, against $-0.19$ and $-0.16$ index by index.

\paragraph{Intermediate spacings.} The per-position analyses use spacings that are powers of two. The grids also include $k = 3$, and repeating the comparison at $k = 2$, 3, and 4 separately gives a mean post-marker change of $-0.18$, $-0.18$, and $-0.17$ over the twenty dependency-task cells at the five open models in Appendix Table~\ref{tab:blindk}, with $k = 3$ matching the sign of its neighbours in 18 of 20. At Qwen2.5-32B we also ran $k = 5$ and 6, which share no divisor with the power-of-two grid: the post-marker change, averaged over the four tasks, is $-0.15$, $-0.14$, $-0.16$, $-0.13$, and $-0.13$ at $k = 2$ through 6, and the within-segment change, averaged the same way, is $+0.04$, $+0.07$, $+0.09$, $+0.09$, and $+0.11$, lower at small $k$ because markers also hurt the within-segment positions of \task{copy\_previous} there. At Qwen2.5-7B, near chance on two of the four tasks, the same means are $-0.09$, $-0.12$, $-0.06$, $-0.05$, and $-0.03$. The loss is a property of the position after a marker, not of a rhythm in the grid.

\paragraph{Other delimiters and longer sequences.} At Qwen2.5-32B we replaced \texttt{|} with a comma or a line break, each with a constraint sentence naming it, and extended \task{running\_xor} and \task{running\_sum\_local} (modulus 4) to $N = 128$ and 256 ($k \in \{4, 16\}$, 20 trials per cell; Appendix Table~\ref{tab:delimiter}). With every delimiter and at every length, the post-marker position is at chance: 0.48 to 0.52 on \task{running\_xor} (chance 0.5) and 0.21 to 0.26 on \task{running\_sum\_local} (chance 0.25). Within-segment positions gain $+0.16$ to $+0.26$ and $+0.09$ to $+0.19$ over the no-marker run at the same positions. Measured against the no-marker run, the post-marker loss shrinks with length (on \task{running\_xor}, $-0.11$ to $-0.14$ at $N = 64$, $-0.03$ to $-0.05$ at 128, and $+0.03$ to $+0.05$ at 256), but only because the no-marker run itself decays toward chance on long sequences; the post-marker position does not improve. The blind spot is a property of the boundary, not of the symbol. The line break gives the largest pooled gain in all six dependency cells, by up to 0.08.

\paragraph{Input markers alone.} Markers in the input with an output asked for without them leave the post-marker position at chance on \task{running\_xor} (Section~\ref{sec:sides}), so the failure follows the segmentation of the input, not the request to emit markers.

\paragraph{Further checks.} Three further checks change nothing: delivering the prompt without the chat template or with other system prompts, raising the modulus of \task{running\_sum\_local} to 16 and 32, and sampling instead of greedy decoding (Appendix~\ref{app:delimiter}, Table~\ref{tab:decoding}).

\section{Repairing the Blind Spot}
\label{sec:repair}

Three interventions target the crossing step directly: changing the rule, changing the instruction, and finetuning.

\paragraph{Reset variants.}
\label{sec:reset}
If the post-marker position fails because it needs the input before the marker, removing that need should repair exactly that position and leave the others alone. We rebuild each dependency task as a \emph{reset} variant: the position right after each marker is treated like the first position of the sequence, so the recurrence restarts at every marker. The prompt states the new rule with a worked example (Appendix~\ref{app:prompts}). Reset variants were run at Qwen2.5-7B and 32B, Qwen3.8-27B, and Gemma-4-31B.

Appendix Figure~\ref{fig:reset} confirms the prediction. On \task{running\_xor}, \task{running\_sum\_local}, and \task{sliding\_window\_max} at the three larger models, the post-marker position rises from 0.28--0.61 to 0.92--1.00, while within-segment positions change by $-0.09$ to $+0.07$. Qwen2.5-7B moves the same way with smaller gains (\task{running\_xor} 0.41 to 0.62; \task{running\_sum\_local} 0.22 to 0.49). Two caveats bound what this shows. The reset rule makes the post-marker position easier by design (on \task{running\_xor} it becomes ``copy this digit''), so the experiment localizes the harm rather than showing a new benefit. And the reset prompt alone shifts no-marker accuracy by $-0.08$ to $+0.06$, so pooled reset effects compare different baselines and are not reported; the per-position comparison uses marker runs only.

\task{copy\_previous} behaves differently. Its post-marker position is also repaired (0.52--0.66 to 0.83--1.00), but its within-segment positions fall at Qwen3.8 ($-0.03$) and Gemma 4 ($-0.07$), and in the original task they are below the no-marker run at eight of the nine models including both Gemini versions, with Mistral level (Appendix Table~\ref{tab:aligned}). Because the output is a shifted copy of the input, an off-by-one slip at a marker can carry through the rest of the segment; we read this as desynchronization, a failure distinct from the post-marker loss. Edit-distance alignment supports this reading: once slips are repaired, \task{copy\_previous}'s within-segment loss shrinks from $-0.24$ to $-0.05$ at Qwen3.8 and from $-0.10$ to $-0.05$ at Qwen2.5-14B, while its post-marker loss remains (Appendix Table~\ref{tab:aligned}).

\paragraph{Instruction wording.} Our instruction calls the bars ``structural boundaries'', which might invite the model to treat each segment as a fresh start. We reran the four dependency tasks at Qwen2.5-32B, Qwen3.8, and Gemma 4 with three other instructions (Table~\ref{tab:wording}): a \emph{neutral} one that only asks for each bar to be copied through; an \emph{across} one stating that the bars are only visual separators and that the position right after a bar uses the input right before it; and the across wording plus two worked examples with bars ($N = 8$, $k = 2$ and 4, produced by the task's own rule), the first annotated at the position right after its first bar. The neutral wording leaves the post-marker loss within 0.03 of the original except at Gemma 4 on \task{copy\_previous} ($-0.29$ to $-0.11$) and \task{sliding\_window\_max} ($-0.07$ to $0.00$). The across wording narrows it in 11 of 12 cells, most at Gemma 4, where it removes it on \task{copy\_previous} ($+0.03$) and \task{sliding\_window\_max} ($+0.09$) and narrows it on \task{running\_sum\_local} from $-0.41$ to $-0.15$, and least at Qwen2.5-32B (by 0.00 to 0.05). The worked examples add to it in 9 of 12 cells: at Gemma 4 most of the repair comes from the wording, while at Qwen2.5-32B the examples add about as much as the wording does. The post-marker position stays below the no-marker run under every wording in 9 of the 12 cells, and at Qwen2.5-32B it remains near chance on \task{running\_xor} (0.57, against 0.69 without markers at the same positions). Within-segment gains move by at most 0.06, except on \task{copy\_previous} at Qwen3.8, where the across wording and examples lift those positions from $-0.24$ to $-0.02$, and at Gemma 4 ($-0.07$ to $+0.01$ under the across wording); no-marker baselines move by at most 0.06. At Qwen2.5-7B, near chance on these tasks, the worked prompt changes nothing except on \task{sliding\_window\_max} ($-0.11$ to $+0.02$). The framing does not create the blind spot, and naming the dependency, with worked examples of the crossing step, removes most of it at Gemma 4, about half at Qwen3.8, and about a third at Qwen2.5-32B: how far a model can be talked out of it depends on the model.

\begin{table}[t]
\centering
\small
\caption[Instruction wording and the blind spot.]{Instruction wording and the blind spot: per-position change from adding markers (marker run minus the no-marker run of the same wording at the same positions). Orig.: ``structural boundaries''; Neutral: copy each bar through; Across: the bars are only visual separators and the rule applies across them; + worked: the across wording with two worked examples with bars, the first annotated at the position after its first bar. Four dependency tasks, $N \le 64$, $k \in \{2,4,8,16\}$, 20 trials per cell.}
\label{tab:wording}
\adjustbox{max width=\linewidth,max totalheight=0.86\textheight}{\input{tables/wording.tex}}
\end{table}

\paragraph{A finetuned separator.} Training can do what prompting only partly does. LoRA adapters trained on \task{running\_xor} and \task{running\_sum\_local} with marked answers as targets remove the blind spot entirely on the trained tasks at Qwen2.5-7B and 32B (post-marker accuracy 1.00), while control adapters trained on the same tasks without markers widen it ($-0.39$ to $-0.64$). On held-out dependency tasks, controls trained without markers for the same number of steps separate what is specific to marker supervision. At 32B the repair of \task{sliding\_window\_max} is ($+0.02$, against $-0.13$ released and $-0.23$ for the step-matched control), but that of \task{copy\_previous} is not, since the step-matched control also removes it ($0.00$, against $-0.16$ released and $+0.05$ for the marked adapter); at 7B the step-matched control repairs \task{copy\_previous} about as well as the marked adapter ($0.00$ and $+0.04$, from $-0.11$) and alone repairs \task{sliding\_window\_max} ($+0.06$, against $-0.07$). The marked adapter's loss of accuracy without markers on held-out tasks (at 32B, 0.86 to 0.75 on \task{copy\_previous} and 0.87 to 0.75 on \task{random\_permutation}) appears almost as much in the step-matched control (0.80 and 0.77), so it comes mostly from the amount of finetuning (Appendix~\ref{app:base}, Table~\ref{tab:lora}). Marker supervision removes the blind spot on the tasks it was shown; on held-out tasks, the only repair specific to it is \task{sliding\_window\_max} at 32B.

\paragraph{Echoing the needed input.} If the post-marker position fails because the input it needs lies across the marker, writing that input again right after the marker might repair it. On the same inputs as the plain marker runs at Qwen2.5-32B, we repeated the digit before each marker in brackets after it (\texttt{0110|[0]1011}), with an instruction that the bracketed digit is a reminder, not part of the sequence. It does not repair the post-marker position (0.53 to 0.55 on \task{running\_xor}, 0.30 to 0.31 on \task{running\_sum\_local}) and lowers the within-segment positions (0.80 to 0.69 and 0.58 to 0.45); 14\% to 24\% of outputs also copy bracketed digits, which scoring removes. Making the needed input visible at the boundary is not enough.

\section{Prefix-State and External Tasks}

On tasks we did not design (Appendix Tables~\ref{tab:external} and~\ref{tab:gemini}), the pattern holds. On CLRS-Text's running-minimum algorithm \citep{markeeva2024clrstext}, where every position depends on the whole prefix, markers hurt wherever the model can do the task: $-0.20$ ($z = 4.5$) at Qwen2.5-32B and $-0.34$ ($z = 7.3$) on Gemini 3.5. On ProcBench's per-character substitution task \citep{fujisawa2024procbench}, a position-local task, markers help ($+0.10$, $z = 6.7$ at Qwen2.5-7B; $+0.20$, $z = 10.8$ at 32B); its strings mix letters and digits, so part of this gain may be tokenization. On string copying \citep{jelassi2024repeat}, where accuracy without markers is above 0.97, the effect is below the significance threshold ($z \le 2.0$). Tagging each letter as a vowel or consonant from pretrained knowledge shows no benefit at Qwen2.5-32B ($+0.01$, $z = 0.5$).

Two further tasks we built with natural readings, the nesting depth of a parenthesis string and a running checksum, carry the whole prefix as state. At Qwen2.5-32B and Gemma 4 markers hurt both, and the within-segment positions fall with the post-marker position instead of gaining (Appendix~\ref{app:external}). The blind spot is confined to one position only when one position is all that needs what lies before the marker.

% =====================================================================
\chapter{Real Sequences}
\label{sec:real}

Every input so far is drawn uniformly at random. Real sequences repeat symbols, have skewed frequencies, and contain words and motifs seen in pretraining. We therefore repeat the design on windows of two public corpora: human promoter DNA from the Nucleotide Transformer benchmark \citep{dallatorre2025nucleotide} (1,584 sequences of 300 bases, test split of its promoter task) and the WikiText-2 test split \citep{merity2017wikitext} reduced to lowercase letters (1,850 paragraphs). Each trial draws one window of $N \in \{16, 32, 64\}$ characters at a random offset, in the contiguous and the spaced format, with 20 trials per cell at Qwen2.5-7B, Qwen2.5-32B, and Gemma-4-31B. Each source has a position-local rule (the complementary base; each letter shifted by one), two one-step dependencies (whether a base is in the same chemical class as the previous one, and copying the previous base; whether a letter comes later in the alphabet than the previous one, and copying the previous letter), and a prefix-state rule (a running count of G and C bases, or of vowels, modulo 4).

In the contiguous format, four of these rules, the same-class, later-letter, and both counting rules, are near or below chance without markers at every model (0.04 to 0.52), so they cannot show a gain or a loss there; the two copy-previous rules were added after seeing this, as the simplest dependency the models can do on these inputs. With spaced characters, three of those cells rise clearly above chance, and they show the synthetic pattern: on the same-class rule at Gemma 4 (0.67 without markers) and Qwen2.5-32B (0.63), and on the later-letter rule at Gemma 4 (0.63), the post-marker position loses $0.22$, $0.13$, and $0.13$ while the within-segment positions gain $+0.06$, $+0.03$, and $+0.05$. Gemma 4 is also above chance on the two spaced running counts (0.71 and 0.50, chance 0.25), where markers change the post-marker position by $+0.02$ and $-0.03$ (Appendix Table~\ref{tab:realfull}).

On the other four (Table~\ref{tab:real}), the synthetic results reappear. The position-local rules gain from markers at every model in the contiguous format, by $+0.12$ to $+0.21$ on DNA and $+0.15$ to $+0.26$ on text. As on letter strings (Chapter~\ref{sec:tokenization}), spacing the bases alone does more than markers do ($+0.17$ to $+0.22$), the post-marker advantage on contiguous DNA at the Qwen models ($+0.27$ and $+0.28$, against $+0.09$ and $+0.17$ within segments) disappears with spaced bases, and with spaced bases markers still help the two Qwen models ($+0.10$ and $+0.13$), while Gemma 4 is at 0.99 without them. The copy-previous rules lose at every model and in both formats, by $0.06$ to $0.39$, including where the model is near ceiling without markers (Qwen2.5-32B: 0.95 on DNA, 0.91 on text), and the post-marker position loses most. Gemma 4 restarts the rule at each marker: on DNA it writes the start symbol N, which the rule gives only at the first position, after 55\% of markers in the contiguous format and 90\% in the spaced format, and at the same positions in fewer than 0.1\% of no-marker outputs, so its post-marker accuracy falls to 0.17 and 0.05; on spaced text it writes the start symbol \# after 53\% of markers. The Qwen models restart after at most 13\% of markers; most of their post-marker errors are out of alignment, and their within-segment positions fall with the post-marker position, as on synthetic \task{copy\_previous}, where alignment-robust scoring attributes that loss to slips at the marker (Section~\ref{sec:reset}). The copy-previous cells therefore combine the blind spot with that desynchronization and show no within-segment gain; the spaced one-step rules above, where the models are above chance, show the gain alongside the post-marker loss, as the synthetic dependency tasks do. Keeping WikiText's spaces and punctuation, to ask whether natural delimiters already act like markers, makes the copy rule too hard for both models we tried (0.17 and 0.31, against 0.86 to 0.88 on the letters alone), so that question stays open. On real DNA and text, then, a position-local rule gains from markers and a dependency that crosses one fails after it.

\begin{table}[!htb]
\centering
\small
\caption[Real sequences.]{Real sequences: windows of human promoter DNA and of WikiText-2 letters, $N \in \{16, 32, 64\}$, 20 trials per cell. No-marker accuracy, pooled marker effect (no-marker mean length-matched, $k = 3$ excluded), and per-position effects at post-marker and within-segment positions against the no-marker run at the same positions ($k \in \{2, 4, 8, 16\}$), in the contiguous and spaced formats. The other four rules are near chance without markers (Appendix Table~\ref{tab:realfull}).}
\label{tab:real}
\adjustbox{max width=\linewidth,max totalheight=0.86\textheight}{\input{tables/real_data.tex}}
\end{table}

% =====================================================================
\chapter{Mechanistic Analysis}
\label{sec:mechanism}

The interventions here are narrower than the rest of this report: mostly at Qwen2.5-32B, head ablations mainly on \task{random\_permutation}, and no interventions on MLPs or paths. They bound what marker-attending heads contribute and locate where the blind spot acts on the prediction; they are not a full account.

\paragraph{Working hypothesis.} To produce output position $i$, a model must find input position $i$ (and, on dependency tasks, position $i - 1$) among many near-identical tokens. We hypothesize that markers turn this global alignment into a local one: with a marker every $k$ characters, the model can keep its place by the markers it has passed and its offset since the last one. This would explain a gain that appears only on long sequences, is spread over the whole segment, depends little on $k$, and shrinks when alignment slips are forgiven in scoring. The same local frame would explain the blind spot: if ``the previous input'' is looked up relative to the current segment, the first position of a segment finds a marker where its predecessor should be. This is a conjecture, not a result: the experiments below test only its simplest version, in which a few marker-attending heads carry the effects, and they rule that version out.

\section{Attention-Head Ablations}

\paragraph{Marker-attending heads.} In Qwen2.5, individual heads attend sharply from each position to the nearest marker (Appendix Figure~\ref{fig:mechanism}b). Heads are ranked by the mean attention mass that all query positions place on the marker tokens, computed per (layer, head) over five held-out trials at the same task, length, and spacing and re-identified for every such cell; the top $n$ are taken across all layers. For the ten heads with the most attention to markers at Qwen2.5-32B, the mean percentile on a previous-token test is 95 and on an induction-head test 31; since heads selected for attending to an adjacent token will score high on a previous-token test, we read this only as ``these are local heads''. The ten lie in layers 1 to 37 of 64, so the selection is not confined to one depth; heads are ranked by attention to markers rather than by causal effect, and we did not search for the heads with the largest marginal effect, so the layer profile of Section~\ref{sec:patching} is the only layerwise localization we offer. The total attention mass on markers rises with density, but so does the number of markers, so we do not claim markers are attention sinks in the sense of drawing more attention than content tokens on average across query positions; at the step right after a marker they do draw more attention than the input that step needs (Section~\ref{sec:patching}).

\begin{table}[t]
\centering
\small
\caption[Head ablation at Qwen2.5-32B (\task{random\_permutation}, $A = 8$, 20 trials per cell, pooled over $k \in \{2, 8\}$).]{Head ablation at Qwen2.5-32B (\task{random\_permutation}, $A = 8$, 20 trials per cell, pooled over $k \in \{2, 8\}$). ``Top 10'': the ten most marker-attending heads, chosen on held-out trials, are blocked from attending to markers (weights zeroed, rows renormalized). Controls: ten random heads blocked from markers, and the top ten blocked from as many random non-marker input tokens (sham). Share of gain removed is $(\text{Markers} - \text{blocked}) / (\text{Markers} - \text{No markers})$; ``Controls'' gives the larger of the two control shares.}
\label{tab:ablation}
\adjustbox{max width=\linewidth,max totalheight=0.86\textheight}{\input{tables/ablation_32b.tex}}
\end{table}

\paragraph{Blocking the ten most marker-attending heads.} We block the ten most marker-attending heads from the markers during generation and compare with two controls (Table~\ref{tab:ablation}). The block acts after the softmax: in each blocked head, the attention weights from every query position to the marker tokens are set to zero and each row is renormalized, so the head's remaining attention is redistributed over non-marker tokens and nothing else in the forward pass changes; a self-test at the start of every run confirms that the blocked weights are exactly zero. Random heads are drawn uniformly from all (layer, head) pairs outside the top set, not matched by layer, though their median layer is close to the top heads' (34 against 30 for ten heads). On digits at Qwen2.5-32B, where markers raise accuracy from 0.41 to 0.96 at $N = 64$, blocking removes 1\% of the gain (paired drop $+0.005 \pm 0.009$). At $N = 64$ it removes 13\% on spaced letters ($+0.065 \pm 0.058$) and 19\% on contiguous letters ($+0.049 \pm 0.030$), against at most 5\% for either control; at $N = 32$ it removes 7\% or less in every format. On spaced letters part of this is one trial ($N = 64$, $k = 8$): in it the 10-head runs and the 50-head controls all close the first segment one character early and then copy the rest correctly but shifted by one, which index scoring counts as 54 errors, while the unblocked run, the top-50 run, and every 200-head run get it right. Without that trial, blocking the ten heads removes 9\% on spaced letters and the controls at most 1\%. The heads that attend most to markers thus carry about none of the gain on digits, a modest and uncertain share on spaced letters, and up to a fifth on contiguous letters. At Qwen2.5-7B, marker gains in the new experiment are small and the drops from blocking are within the range of the controls (for example, $+0.03 \pm 0.05$ on spaced letters at $N = 32$, against $+0.03$ for random heads). Blocking ten of the model's 2,560 heads is a narrow intervention, and the benefit may be distributed; what we can say is that it is not concentrated in the heads that attend to markers most.

\paragraph{Head-count dose response.} Blocking the 50 most marker-attending heads removes 1\% to 4\% of the gain and the top 200 (8\% of the model's heads) about a quarter, 28\% on spaced letters and 21\% on digits, while 200 random heads blocked from the markers, or the top 200 blocked from as many random non-marker tokens, remove nothing (Appendix Figure~\ref{fig:mechanism}c). The benefit runs partly through attention to the markers, but thinly spread: even 200 heads carry about a quarter of it.

\paragraph{Caveats of the cross-run comparison.} Two features of the comparison qualify it: in the marked run the step that predicts the answer after a marker is read at the output bar rather than at the previous answer digit, and every patched token sits at a later absolute position than its clean counterpart, so a single patched token carries a rotary phase that does not match its neighbours, which patching every input avoids. Without markers the input before and after the marker position are adjacent, so early local heads can already have copied the earlier one into the later one's representation, a step the bar between them interrupts; this is a further reason the input after the marker is the main patching site.

\paragraph{Attention to the needed input.} On \task{copy\_previous} at $N=32$, a marker after every character keeps the marker pattern aligned with the one-step lookback, and attention from each output position to the input it needs is high (0.535). At $k=2$, every other lookback crosses a marker; attention to the needed token falls to 0.101 and accuracy drops from 0.88 to 0.65, below the no-marker baseline of 0.83. At $k=8$ both recover (Appendix Figure~\ref{fig:mechanism}d). The attention here is the mean over the ten heads, chosen separately at each spacing, that attend most to the needed input, averaged over all output positions; part of the recovery at $k = 8$ is therefore that only one position in eight follows a marker, against every other position at $k = 2$, and Section~\ref{sec:patching} isolates the post-marker step. This is the attention-level counterpart of the blind spot.

\paragraph{A hybrid-attention model.} Qwen3.8-27B, in which 48 of 64 layers use Gated DeltaNet linear attention \citep{yang2025gated} rather than softmax attention, shows the tokenization-independent gain and the blind spot comparably to the other strong models.

\FloatBarrier
\section{Input Versus Output Markers}
\label{sec:sides}

The intervention changes two things at once: it segments the input, and it asks for a segmented
output. To separate them we ran a $2 \times 2$ at Qwen2.5-32B on \task{running\_xor} and
\task{random\_permutation} with digits ($N \in \{32, 64\}$, $k \in \{4, 16\}$, 20 trials),
varying whether the input carries markers and whether the output is asked to. The three
marker-bearing conditions are built from the same underlying string per trial, so they are
paired.

Asking for the format through the prompt alone is confounded by partial compliance: asked to omit the markers from its output, the model still emits them in 28\% to 62\% of generations, and asked to insert them into an unmarked output, it often does not (on \task{random\_permutation} it never produced the requested placement). The interaction there is already positive ($+0.14$ and $+0.16$, both intervals excluding zero), but the clean test constrains the decoder.

We reran the same four cells with a logits processor that
counts the answer characters as they are generated and either forces the marker token when one is due or removes every marker-bearing token from consideration, so that the output format is constrained by the decoder rather than requested solely through the prompt (Table~\ref{tab:sidesforced}). The decoder is greedy and the mask only removes tokens, so it cannot change the scores of the tokens it leaves (Appendix~\ref{app:delimiter} gives the procedure). The constraint holds exactly in the
no-marker, input-only, and both-sides cells; in the output-only cell the whole output has the requested format in 64\% to 75\% of
generations on these two tasks, but within the first $N$ answer characters, the part that is scored, it is exact in every generation: the shortfall is markers written after the answer, while the model keeps generating. Constraining does
not itself move the results: the no-marker and both-sides cells reproduce their unconstrained
counterparts to within 0.03.

Under constrained decoding, on these two tasks at Qwen2.5-32B, \emph{neither side alone helps, and both hurt}. Markers in the input
with an output genuinely free of them give $-0.03$ on \task{running\_xor} and $-0.15$ on
\task{random\_permutation}; markers forced into the output of an unmarked input give $-0.05$ and
$-0.16$. Only the combination helps, by $+0.10$ and $+0.39$. The interaction is $+0.19$ (95\%
interval $[0.15, 0.24]$) and $+0.70$ ($[0.61, 0.79]$). The input-only cell under enforcement scores 0.40 on \task{random\_permutation}, close to
the 0.42 of the unenforced generations that happened to comply. On the tasks and model tested, a marker helps only when it appears on both sides, where it can align the output being written with the input being read; segmenting the input while writing an unsegmented output is worse than not segmenting at all. Restricted to the output-only generations whose whole output has the requested format, a selected and easier subset, accuracy is 0.62 on \task{running\_xor} and 0.45 on \task{random\_permutation}, still below the no-marker cell, and the interaction stays positive ($+0.18$ and $+0.64$). Two more dependency tasks at Qwen2.5-32B, \task{copy\_previous} and \task{running\_sum\_local} ($M = 4$), repeat the pattern (Table~\ref{tab:sidesforced}): input alone $-0.12$ and $-0.04$, output alone $-0.06$ and $-0.05$, both sides $+0.02$ and $+0.08$, with interactions of $+0.20$ ($[0.12, 0.27]$) and $+0.17$ ($[0.13, 0.22]$), so the interaction holds on all four tasks tested at this model. On these two tasks the whole output has the output-only format in only 32\% to 35\% of generations, because the model runs past the required length and keeps inserting markers, but the scored first $N$ characters are exact in every generation (Table~\ref{tab:sidesforced}). The fully compliant subset scores 0.80 on \task{copy\_previous} and 0.41 on \task{running\_sum\_local}, 0.03 and 0.02 above the no-marker cell, and the interaction computed with it stays positive ($+0.11$ on both).

Fourth, the blind spot follows the input markers alone. With markers in the input and an output
enforced to carry none, the position right after a marker is at chance on \task{running\_xor}
(0.53, chance 0.5), as it is when markers appear on both sides (0.51); on \task{copy\_previous} and \task{running\_sum\_local} the input-only and both-sides cells agree to within 0.02 at that position as well. Segmenting the input is
what breaks the cross-boundary dependency; emitting markers is not.

\begin{table}[!htb]
\centering
\small
\caption[The same 2$\times$2 with the output format constrained by a logits processor rather than requested solely through the prompt (Qwen2.5-32B.]{The same 2$\times$2 with the output format constrained by a logits processor rather than requested solely through the prompt (Qwen2.5-32B; $N \in \{32, 64\}$, $k \in \{4, 16\}$ pooled, 20 trials per cell;) Format realized: the share of generations whose whole output has the requested markers (full output) and whose first $N$ answer characters, the scored part, do (first $N$); acc.\ if so: accuracy over the generations whose first $N$ characters comply. Full-output shortfalls in the output-only cells are markers written after the answer, where the model keeps generating. Change is against the no-marker cell, and the interaction is $A_{11} - A_{10} - A_{01} + A_{00}$, computed from unrounded cell means, with a bootstrap interval. Post-marker is accuracy at the position right after a marker.}
\label{tab:sidesforced}
\adjustbox{max width=\linewidth,max totalheight=0.86\textheight}{\input{tables/marker_sides_forced.tex}}
\end{table}

\paragraph{Two 7B-class models.} At Qwen2.5-7B and Mistral-7B the constrained $2\times2$ is near chance in six of its eight task--model blocks; where a model is above chance the interaction holds ($+0.16$ on \task{random\_permutation} at Qwen2.5-7B, $+0.20$ on \task{running\_xor} at Mistral).

% =====================================================================
\section{Activation Patching and Cross-Boundary Attention}
\label{sec:patching}

Where is the blind spot represented? We teacher-force Qwen2.5-32B on \task{running\_xor} and \task{running\_sum\_local} ($M = 4$; $N = 32$, $k = 4$, 20 inputs per task) with and without markers, read the log-odds of the correct digit against the reset digit at each post-marker step where the two differ (69 of 140 such positions on \task{running\_xor}, 100 of 140 on \task{running\_sum\_local}; averaged within an input, then over inputs), and then overwrite the marked run's residual stream at one layer, or at every layer, with the no-marker run's at chosen positions (Appendix~\ref{app:patching}, Table~\ref{tab:patching}). Without patching the marked run prefers the reset digit by 6 to 7 nats; the gap between the clean and the marked run's log-odds, the denominator of every recovery ratio, is 7 to 9 nats. Patching every input digit recovers 75\% and 86\% of the gap, the input digit right after the marker alone 59\% and 64\%, the digit before the marker 12\% and 1\%, and the answer digit already written 18\% and 20\%. Patched one layer at a time, the recovery from the input digits peaks at 81\% and 94\% around layer 34 of 64 and is gone by layer 50. Patching thus identifies the residual-stream representation of the first input digit after a marker as the most influential site for the cross-boundary failure, with the strongest single-layer recovery around layer 34 of 64 and none after layer 50; because the clean and marked runs differ in token positions and context, it shows where the difference acts, not where it originates. The same procedure at Gemma 4 (60 layers) and Qwen3.8 (64 layers, mostly linear attention) gives the same ordering of sites: every input digit recovers 60\% and 61\% of the gap at Gemma 4 and 46\% and 74\% at Qwen3.8, the input digit after the marker alone 48\% and 30\% at Gemma 4 and 48\% and 69\% at Qwen3.8, the digit before the marker at most 6\% at Gemma 4 and 21\% at Qwen3.8, and the answer already written 15\% to 20\% at Gemma 4 and 8\% to 35\% at Qwen3.8; the single-layer peak of the every-input curves is around layer 26 of 60 at Gemma 4 and layers 25 to 34 of 64 at Qwen3.8, with nothing after layer 50 at either (Appendix~\ref{app:patching}). Qwen3.8's clean--marked gap is 2.2 nats, so its ratios are the least precise of the three. At Qwen2.5-7B, at chance on these tasks at $N = 32$ even without markers, the gap is too small to patch (Appendix~\ref{app:patching}).

\paragraph{Attention to the input before the marker.} Patching shows which representation the failure acts through; the attention at the failing step shows what the model reads. In the same teacher-forced runs at Qwen2.5-32B we record, for the step that writes each answer character, the attention placed on the input character before it, per layer and head (Appendix~\ref{app:patching}, Table~\ref{tab:attention}). Without markers, a band of heads in layers 35 to 45 looks back at that character, the strongest of them with 0.22 to 0.60 of its attention (Figure~\ref{fig:mechlayers}, bottom). With markers the band is intact at within-segment steps, which keep 89\% to 96\% of the unmarked lookback in layers 35 to 45, and gone at the step after a marker, which keeps 30\% to 38\% of it on average over heads and 23\% to 27\% for the strongest head; in layers 35 to 45 the strongest head's attention to the needed input falls to 0.04 to 0.07. At that step the marker receives 0.56 to 0.59 of the attention, summed over the 64 layers (so out of a possible 64) and averaged over heads, more than the 0.44 to 0.46 the needed input receives without markers. This is why patching the input before the marker recovers little: the step after a marker barely reads it, while its attention to the input after the marker stays high (0.67 and 0.98; Appendix~\ref{app:patching} discusses the comparison's caveats). Single-layer patching recovers the prediction only up to the layers of this band (Figure~\ref{fig:mechlayers}, top): once the lookback has run, replacing the input's representation no longer helps. The same collapse appears at Gemma 4 and Qwen3.8, where attention from the step after a marker to the input before it falls from 0.47 to 0.57 in the clean run to 0.08 to 0.13 (Table~\ref{tab:attention}). In these runs, however, the step after a marker is read at the output bar, so the bar is also the query. With markers in the input only, the query is the previous answer digit, and the lookback largely survives: the input before the marker keeps 0.37 to 0.40 of the attention at Qwen2.5-32B (against 0.44 to 0.46 clean), 0.30 to 0.33 at Gemma 4, and 0.40 to 0.42 at Qwen3.8, while the input at the step's own position falls to 0.19 to 0.35 (against 0.59 to 0.78 clean) and the marker takes 0.40 to 0.52. Yet with input markers only that position still fails (Section~\ref{sec:sides}). The collapse of the lookback is therefore tied to the output bar, and the input-only failure goes with attention drawn from the current input to the marker rather than with a lost lookback; neither measurement shows which carries the dependency.

\begin{figure}[!t]
\centering
\includegraphics[width=\linewidth]{figures/fig_mechanism_layers.pdf}
\caption[Where the blind spot acts at Qwen2.5-32B ($N = 32$, $k = 4$, 20 teacher-forced inputs per task), on a shared layer axis.]{Where the blind spot acts at Qwen2.5-32B ($N = 32$, $k = 4$, 20 teacher-forced inputs per task), on a shared layer axis. Top: log-odds recovery when the unmarked run's residual stream replaces the marked run's at one layer, by patched site (0 is the marked run, 1 the unmarked run). Bottom: attention from the step that writes each answer character to the input character before it, strongest head per layer, without markers at the post-marker indices and with markers at within-segment and post-marker steps. The lookback band in layers 35 to 45 is gone at the step after a marker, and patching recovers the prediction only in the layers up to that band.}
\label{fig:mechlayers}
\end{figure}

\chapter{Discussion and Conclusion}
\label{sec:practical}
\label{sec:limitations}
\label{sec:conclusion}

The experiments separate three things a content-free marker does. On letter strings it mostly re-tokenizes; with one token per character it still helps the stronger models keep their place on long sequences; and with a marker every two or more characters, the position right after it fails whenever its output depends on the input before the marker, while rewriting the rule so that position needs nothing before the marker repairs exactly that position. Table~\ref{tab:summary} collects the main effect sizes.

\begin{table}[!htbp]
\centering
\small
\caption[Summary of the main effects.]{Summary of the main effects. Pooled effects are over $N \le 64$ with the no-marker mean length-matched; per-position effects are against the no-marker run at the same positions. ``Strong models'': Qwen2.5-14B and 32B, Qwen3.8-27B, Gemma-4-31B.}
\label{tab:summary}
\adjustbox{max width=\linewidth,max totalheight=0.86\textheight}{\begin{tabular}{p{0.43\linewidth} p{0.30\linewidth} p{0.19\linewidth}}
\toprule
Effect & Size & Where \\
\midrule
\multicolumn{3}{l}{\emph{1. On letter strings, markers mostly re-tokenize}} \\
Markers on contiguous letters, pooled & $+0.06$ to $+0.26$ across models & Chapter~\ref{sec:tokenization} \\
Spacing the letters instead, no markers, at Qwen and Gemma & $+0.16$ to $+0.33$ (beats markers) & Figure~\ref{fig:formats} \\
Post-marker advantage on letters, gone with spaced letters or digits & $+0.07$ to $+0.34$ vs within 0.04 & Figure~\ref{fig:position} \\
\midrule
\multicolumn{3}{l}{\emph{2. The gain survives one token per character}} \\
Share of remaining errors removed, strong models, spaced or digits & 33\% to 83\% (23 of 24 cells significant) & Chapter~\ref{sec:segmentation} \\
Gain at $N = 64$, strong models, all 24 control cells & $+0.07$ to $+0.37$, all positive & Appendix~\ref{app:stats} \\
Persists to $N = 256$ at Qwen2.5-32B, any of three delimiters & 56\% of remaining errors removed & Table~\ref{tab:delimiter} \\
Blocking 10 / 200 marker-attending heads (spaced letters or digits) & removes $\le 13\%$ / 21--28\% & Table~\ref{tab:ablation}, Figure~\ref{fig:mechanism}c \\
\midrule
\multicolumn{3}{l}{\emph{3. The post-marker blind spot on dependency tasks}} \\
Post-marker position vs no-marker run, strong models & $-0.15$ to $-0.26$ (Gemini $-0.34$, $-0.43$); at chance & Figure~\ref{fig:position} \\
Next position & recovers ($+0.20$ on \task{running\_xor}) & Chapter~\ref{sec:blindspot} \\
Reset variants (dependency severed at the marker) & post-marker rises to 0.92--1.00 & Appendix Figure~\ref{fig:reset} \\
Survives alignment-robust scoring, commas and line breaks, $N = 256$, modulus 32, sampling, other prompt deliveries & post-marker stays at chance & Chapter~\ref{sec:blindspot} \\
Across instruction naming the dependency & narrows it in 11 of 12 cells; removes it only at Gemma 4 on two tasks & Table~\ref{tab:wording} \\
Across wording plus two worked examples & removes most of the loss at Gemma 4, half at Qwen3.8, a third at Qwen2.5-32B; mostly the wording at Gemma 4, wording and examples alike at Qwen2.5-32B & Table~\ref{tab:wording} \\
LoRA on marked targets (Qwen2.5-7B, 32B) & removes the loss on the trained tasks, where unmarked controls widen it; against step-matched controls, the only held-out repair specific to it is one task at 32B; its cost in no-marker accuracy comes mostly from the amount of training & Table~\ref{tab:lora} \\
Input markers alone (output constrained marker-free) & post-marker still at chance (0.53) & Table~\ref{tab:sidesforced} \\
Patching and attention at the post-marker step & input after the marker is the main site; lookback attention collapses when the output carries markers; with input markers only the current input loses attention to the marker & Section~\ref{sec:patching} \\
\midrule
\multicolumn{3}{l}{\emph{Real sequences (promoter DNA, WikiText-2 letters; three models)}} \\
Position-local rules, contiguous & $+0.12$ to $+0.26$ & Table~\ref{tab:real} \\
Copy previous, both formats & $-0.06$ to $-0.39$; Gemma 4 restarts after 55--90\% of markers & Chapter~\ref{sec:real} \\
Other one-step rules, spaced, where above chance & post-marker $-0.13$ to $-0.22$, within $+0.03$ to $+0.06$ & Chapter~\ref{sec:real} \\
\midrule
\multicolumn{3}{l}{\emph{Checks on the design}} \\
Cost of the baseline's marker sentence, 300 cells & $+0.006$ $[0.002, 0.009]$ & Appendix~\ref{app:stats} \\
Input/output 2$\times$2 interaction, constrained decoding & $+0.17$ to $+0.70$ on four tasks at Qwen2.5-32B, intervals exclude 0; $+0.16$ and $+0.20$ at 7B-class models above chance & Table~\ref{tab:sidesforced} \\
\bottomrule
\end{tabular}}
\end{table}

\paragraph{Practical guidance.} Table~\ref{tab:decision} summarizes the recommendations. For letter strings, spacing the characters beats markers at every Qwen and Gemma model we test, matches them at Mistral and Phi, and costs fewer tokens: 48 spaced letters are 48 tokens, against 95 with a marker after every letter. On top of spacing, a marker every 4 to 16 characters adds a further gain at the stronger models from $N \approx 32$, provided the model reproduces the markers in its output, since marking either side alone did not help in the constrained-decoding ablation; where a model does not reproduce them reliably (prompted compliance is 26\% to 48\% on the dependency tasks at the 7B-class models, Appendix Table~\ref{tab:compliance}), they can be inserted by constrained decoding or stripped after generation. Sparse markers also stop models from running past the required length: on contiguous-letter tasks at $N = 64$ to 128, 18\% to 95\% of outputs run past it without markers at Qwen2.5-32B, Qwen3.8, and Gemma 4, against 0\% to 13\% with a marker every 8 characters, so on contiguous \task{cyclic\_shift} at Qwen2.5-32B a marker every 4 or 8 characters takes less time than none from $N = 64$ on (0.66 times as long at $N = 64$ and 0.44 at $N = 128$ for $k = 8$), while $k \le 2$ and shorter sequences cost more (Appendix Table~\ref{tab:wallclock}); we did not time the one-token-per-character formats. True divide-and-conquer, one call per segment, holds accuracy nearly flat (0.93 at $N = 8$ to 0.92 at 128 on contiguous \task{cyclic\_shift}) but costs about three times the wall-clock of the marker condition. Where an output depends on the input just before a marker, the marker should not be placed there, or the rule should be restated so that the position after each marker needs nothing before it.

\begin{table}[!htbp]
\centering
\small
\caption[When to space characters, when to add markers, and when to avoid them.]{When to space characters, when to add markers, and when to avoid them. Based on the models and tasks tested here (the limitations discussed below bound its scope).}
\label{tab:decision}
\adjustbox{max width=\linewidth,max totalheight=0.86\textheight}{\begin{tabular}{p{0.45\linewidth} p{0.49\linewidth}}
\toprule
Situation & What to do \\
\midrule
Characters run together and the tokenizer merges them (letter strings) & Put a space between characters first. At the Qwen and Gemma models this beats markers and costs fewer tokens \\
Sequence long enough that accuracy has started to fall (from about $N = 32$ here) & Add a marker every 4 to 16 characters and ask the model to reproduce them in its output; this combination gave the largest gains in the four-task ablation (Section~\ref{sec:sides}). The gain depends little on the spacing, and sparse markers cost little. If the model does not reproduce the markers reliably, insert them by constrained decoding or strip them afterwards \\
Short sequences, where the model is near ceiling & Markers add nothing; skip them \\
Each output depends only on the input at the same position & Markers generally improve accuracy on the position-local tasks tested, particularly at stronger models \\
An output depends on the input just before a marker & Do not place a marker there, or restate the rule so the position after each marker needs nothing before it; if the marker must cross the dependency, say so in the instruction and add worked examples of the crossing step, which together removed most of the loss at Gemma 4, half at Qwen3.8, and a third at Qwen2.5-32B \\
Latency matters & On contiguous letters from $N = 64$ on, a marker every 4 to 8 characters was faster than none, because it stops the model from running past the required length; denser markers and shorter sequences cost more, and one-token-per-character formats were not timed \\
\bottomrule
\end{tabular}}
\end{table}

\paragraph{The marker symbol and its placement.} On contiguous-letter \task{cyclic\_shift} at two models, the choice of symbol matters less than its familiarity: \texttt{;} and \texttt{\#} keep accuracy above the no-marker baseline as \texttt{|} does, while a reserved special token the model has never seen in this role falls far below it, and the benefit reverses at Qwen2.5-32B under a terse cipher-style phrasing of the instruction. Moving every marker by up to one position leaves the blind spot in place on the dependency tasks, and commas and line breaks behave like \texttt{|} there, so neither effect depends on exact periodicity or on this particular symbol (Appendix~\ref{app:marker}).

\paragraph{Mechanism.} The experiments bound the mechanism rather than settle it. Marker-directed attention contributes to the gain, but thinly: blocking 200 marker-attending heads removes about a quarter of it, fewer heads about nothing, and we did not test MLP or path-level interventions. Marking only one side hurts and only marking both helps, while the cross-boundary failure follows the input markers alone. The failure is consistent with the model resolving ``the previous input'' within the current segment: with markers on both sides, the heads that read the previous input stop reading it at the step after a marker at all three models we measured, and patching locates the failure in the representation of the first input after the marker; with input markers only, the lookback survives but the current input loses attention to the marker, so the attention measurements do not isolate a single cause, and two tasks at three models do not establish the mechanism in general. Separators that act as attention anchors help sequence-level questions over long numerical inputs \citep{sun2026sepseq}; for per-position transformations the anchor helps only when the output is segmented with the input, and it is also where the cross-boundary dependency is lost. Combining patching with causal head gating \citep{nam2025headgating}, which names the heads involved, is the natural next probe. Retokenization sampling \citep{jain2026retokenization} finds that behaviour stays sensitive to how the same bytes are segmented deep into the network; our controls show that a segmentation imposed by content-free tokens, with the per-character tokenization fixed, persists in the same way.

\paragraph{Limitations.} The open models are instruction-tuned, five of them 2024-era; base checkpoints of Qwen2.5-7B and 32B show the blind spot at close to its instruct size (Appendix~\ref{app:base}), so instruction tuning enlarges the effects but does not create them. Apart from the real windows of Chapter~\ref{sec:real}, the tasks are synthetic character-level rules on letters or digits, with alphabets or moduli of at most eight symbols in the core grids (32 in one follow-up), dependencies of one or two positions, and lengths up to $N = 64$ (256 in one follow-up); other tokenizer families, non-Latin scripts, and thousands of characters are untested. The open models are decoded greedily; Gemini 3.6 ignores the temperature parameter, Gemini 3.5 has 3 trials per cell, the two were queried three months apart, Gemini's tokenization is known only through API token counts, and at 0.94 to 1.00 without markers on the controls Gemini gives no test of the tokenization-independent gain. Symbol, wording, placement, and delimiter ablations cover few tasks and models. We test only content-free markers and only rotary position embeddings; ALiBi or FIRE encodings, separator tokens such as \texttt{[SEP]}, applied settings such as log lines or tables, and byte-level models are effectively untested (EvaByte-SFT, the one openly available byte-level chat model, copies its input instead of applying the rules; Appendix~\ref{app:base}). On method: position-wise scoring penalizes a dropped character at every later position, which alignment-robust scoring corrects without removing the blind spot (Appendix Table~\ref{tab:aligned}); marker and no-marker runs use different inputs, at a small cost in precision (Appendix Table~\ref{tab:paired}); the no-marker baseline keeps the sentence describing the markers, at a measured cost of $+0.006$ at the four stronger models and $-0.003$ at Qwen2.5-7B; the reset rule makes the post-marker position easier by construction; and per-position claims concern spacings of two or more.

In sum, a content-free marker is a cheap, single-call intervention whose benefit and whose failure mode are consistent with one explanation, that it gives the model a local frame, although the experiments establish the two effects more firmly than how the model uses the boundaries. Separate the characters, mark long sequences sparsely on both input and output, and keep markers away from any position whose output depends on what came just before them.

\paragraph{Reproducibility.} Code, prompts, raw generations, and analysis scripts are at \url{https://github.com/faezehkhazaee97/AvalancheBreak}. All inputs are generated from deterministic seeds, and every prompt can be regenerated from (task, format, length, spacing, alphabet, trial).

\FloatBarrier
\bibliographystyle{tmlr}
\bibliography{refs}

\appendix
\chapter{Supplementary Results}



\section{Prompts}
\label{app:prompts}

\paragraph{\task{cyclic\_shift} ($N = 8$, $A = 4$, $k = 4$).}
\begin{small}
\begin{verbatim}
You are a deterministic sequence transformer.
Task: Cyclic character shift.
Alphabet: A, B, C, D
Rule: Replace every character with the next one in the alphabet.
Mapping: D -> A, and A -> B.
Constraint: The input contains vertical bars (|). These are structural
boundaries. DO NOT shift them. Copy them exactly to the output.
No spaces, no explanations.

Input: CABA|CCDA
Output:                 (expected: DBCB|DDAB)
\end{verbatim}
\end{small}

The mapping line is the literal string the prompt code generates: it shows the wrap-around step and the first step only, and the rule line above it states the full shift.

\paragraph{Spaced and digit formats.} The spaced version of this cell shows the input as \texttt{C A B A | C C D A} and replaces ``No spaces, no explanations.'' with ``Write the output in the same format as the input: every character separated by a single space. No explanations.''; everything else is unchanged. The digit version relabels $A, B, C, D$ as $0, 1, 2, 3$ (input \texttt{2010|2230}), with the task line, rule, and mapping worded for digits (``Cyclic digit shift'', ``Replace every digit with the next one in the alphabet''). The same seed draws the same positions in all three formats.

\paragraph{\task{running\_xor}, original and reset ($N = 8$, $k = 4$).} Both prompts state the rule $y_i = x_i \oplus x_{i-1}$ with the first digit copied. The reset prompt adds: \emph{``Immediately after each boundary, treat that position exactly like the first digit of the whole input---copy it as-is (nothing before a boundary counts as `previous'), then resume the XOR rule from there until the next boundary,''} with the worked example \texttt{1011|0010} $\rightarrow$ \texttt{1110|0011}. For input \texttt{1101|1111}, the original expects \texttt{1011|0000}, while the reset variant expects \texttt{1011|1000}: only the digit after the marker differs.

\paragraph{\task{sliding\_window\_max} ($N = 8$, $k = 4$).} The prompt states the rule with a worked example and adds: \emph{``DO NOT include them [the bars] in the window or count them as digits.''} For input \texttt{5170|1247} the expected output is \texttt{5577|7247}.

\paragraph{More verbatim prompts.} The prompts below are regenerated from the task code with the grid's seeds ($N = 8$, $k = 4$, trial 0), so they are exactly what the models received, except that long lines are wrapped here to fit the page and the arrow and XOR symbols are printed as \texttt{->} and \texttt{\^{}}. For the digit version of \task{conditional\_shift} and for the neutral, across, and worked wordings of \task{running\_xor}, only the lines that differ from the prompt shown before them are printed.

\input{tables/prompts_verbatim.tex}

\section{Trials per cell}
\label{app:trials}

Table~\ref{tab:trials} gives trials per cell (one cell = task, model, length, density, and alphabet or modulus value) for the data pooled in Figures~\ref{fig:formats}, \ref{fig:profile}, and~\ref{fig:split}. Gemini 3.5 Flash uses 3 on the letter and dependency tasks; Gemini 3.6 Flash uses 5 on all 13 task--format variants.

\paragraph{Closed models.} The closed models are two versions of Gemini Flash, queried through the API with the same prompts and seeds; the resolved model name and timestamp are stored with every generation. Gemini 3.5 Flash \citep{gemini35flash} was run on the letter and dependency tasks, but not on the two controls, between 2026-08-11T23:28 and 23:45 UTC at temperature 0 with thinking disabled (\texttt{thinking\_budget} 0), the closest available match to the open models' greedy decoding, under the versioned model name rather than the \texttt{-latest} alias, and every response reported \texttt{gemini-3.5-flash} as the model served; its usage metadata likewise reports a total equal to prompt plus output tokens on all 2,040 calls. On 2026-10-08 Google deprecated \texttt{gemini-3.5-flash} and began routing its requests to \texttt{gemini-3.6-flash} \citep{geminichangelog}, so the controls could not be added at that model; instead we ran the full grid of all 13 task--format variants (the seven tasks with both controls) on Gemini 3.6 Flash \citep{gemini36flash} between 2026-11-12T07:56 and 10:00 UTC with 5 trials per cell. Every request sent \texttt{thinking\_budget} 0 (the runner probes the setting once per model name and reuses it for every call), which this model accepts: on every one of the 4,930 calls the usage metadata reports no thought tokens and a total equal to prompt plus output tokens, and every answer is full length, whereas a request at the model's default thinking level spends about 160 thought tokens and, under our output budget, truncates the answer. Gemini 3.6 Flash deprecates and ignores the temperature parameter, so its outputs are samples at the model's fixed settings rather than greedy decodes: repeating an identical request four times with thinking off gave two to four distinct outputs on each of three cells we checked. Its trials are therefore independent samples as well as independent inputs.

\begin{table}[!htbp]
\centering
\small
\caption[Trials per cell, $N \le 64$.]{Trials per cell, $N \le 64$. Q2.5: Qwen2.5-Instruct; Q3.8: Qwen3.8-27B; Gemma 4: Gemma-4-31B-it. Reset variants at Qwen2.5-7B cover three of the four dependency tasks (\task{sliding\_window\_max} was not run in reset form there).}
\label{tab:trials}
\adjustbox{max width=\linewidth,max totalheight=0.86\textheight}{\begin{tabular}{lccccccc}
\toprule
Task & Q2.5-7B & Q2.5-14B & Q2.5-32B & Mistral & Phi & Q3.8 & Gemma 4 \\
\midrule
\task{cyclic\_shift} & 20 & 20 & 20 & 40 & 20 & 20 & 20 \\
\task{random\_permutation} & 40 & 40 & 40 & 40 & 20 & 20 & 20 \\
\task{conditional\_shift} & 40 & 20 & 40 & 40 & 20 & 20 & 20 \\
\task{copy\_previous} & 40 & 40 & 80 & 40 & 40 & 40 & 40 \\
\task{running\_xor} & 40 & 40 & 40 & 40 & 20 & 40 & 40 \\
\task{running\_sum\_local} & 40 & 40 & 80 & 40 & 20 & 40 & 40 \\
\task{sliding\_window\_max} & 20 & 20 & 40 & 20 & 20 & 20 & 20 \\
Reset variants & 80 & -- & 20 & -- & -- & 20 & 20 \\
\midrule
Spaced versions (each) & 20 & 20 & 20 & 20 & 20 & 20 & 20 \\
Digit versions (each) & 20 & 20 & 20 & 20 & 20 & 20 & 20 \\
\bottomrule
\end{tabular}}
\end{table}

\paragraph{Checkpoints.} Hugging Face revisions: Qwen2.5-7B-Instruct \texttt{a09a35458c70},
Qwen2.5-14B-Instruct \texttt{cf98f3b3bbb4}, Qwen2.5-32B-Instruct \texttt{5ede1c97bbab},
Mistral-7B-Instruct-v0.3 \texttt{c170c708c41d}, Phi-3.5-mini-instruct \texttt{2fe192450127},
Qwen3.8-27B \texttt{37561816e66b}, Gemma-4-31B-it \texttt{842da3794eaa}.

\paragraph{Generation budget.} The no-marker run and every spacing $k \ge 2$ share one budget, $1.5N + 20$ new tokens; $k = 1$ uses $2N + 20$. To check that the budget does not drive any effect, we re-tokenized every stored raw output of the seven open models ($N \le 64$; 86,400 generations) and counted those that used the whole budget while giving fewer than $N$ answer characters on the first line. Raw generations are stored for every spaced-letter and digit control run, every Qwen3.8 and Gemma 4 run, and, at Qwen2.5-32B, the second \task{sliding\_window\_max} batch and the $N = 128$ runs; the 80-trial second batches of \task{copy\_previous} and \task{running\_sum\_local} at 32B (Table~\ref{tab:trials}) do not store it. The original contiguous-letter and dependency grids of the five earlier models (Qwen2.5-7B, 14B, and the first 32B batch, Mistral, and Phi) predate raw-generation logging and retain only the clipped answer, which is why some marker-compliance entries elsewhere are unavailable. That count is an upper bound on cut-off answers, because it also counts short answers followed by unrequested text. The share is 0.58\% without markers, 0.49\% at $k \ge 2$, and 0.31\% at $k = 1$; at the four stronger models it is 0.04\%, 0.02\%, and 0\%. Nearly all cases are Phi-3.5-mini on spaced letters (up to 19\% of one cell without markers, against 14\% with), and on the dependency tasks the share is 0.03\% or less in every class. Because of the logging gap, the dependency-task compliance rows of Appendix~\ref{app:wfr} for Qwen2.5-7B, 14B, and 32B, Mistral, and Phi rest only on the runs that do store raw output: \task{sliding\_window\_max}, added later with logging on, and at Phi also \task{copy\_previous}.

\section{How the uncertainty was computed}
\label{app:stats}

Pooled effects use the length-matched estimator of Chapter~\ref{sec:setup}, with analytic
standard errors over generations. Two checks ask how much that choice matters. Both resample 2{,}000 times,
over the 91 pooled task--model cells of the three position-local families and the dependency family at the seven open models. The 125 cells of the main analysis also include the 15 reset cells and the 19 Gemini cells (6 at Gemini 3.5 with 3 trials each, 13 at Gemini 3.6 with 5), which are not resampled.

\paragraph{Resampling generations reproduces the analytic result.} Bootstrapping generations
within each condition and length gives a median $|z|$ within 3\% of the analytic value, and the
Bonferroni verdict ($|z| \ge 3.54$) changes in 1 of 91 cells. The analytic standard errors are
therefore not the weak point.

\paragraph{Resampling whole cells is far wider.} Treating each (length, spacing, alphabet) cell
as the unit instead, so that whole cells enter or leave a resample, widens the median interval
from 0.046 to 0.155 and leaves 16 of 91 cells above the Bonferroni threshold instead of 62. Two
things drive this, and neither is a defect in the data. First, a pooled cell deliberately spans
lengths at which the effect is absent: at $N \le 16$ these models are near ceiling and markers
do nothing, which is itself Finding~2. Averaging those in and then asking how much the average
moves when lengths are resampled mostly measures that by-design heterogeneity. Second, a pooled
cell contains only 10 to 25 (length, spacing, alphabet) cells, and a bootstrap over so few
clusters has little power.

Restricting to the length where the claim is actually made removes the first problem. At
$N = 64$, all 24 spaced-letter and digit cells at the four stronger models have a positive
effect, between $+0.07$ and $+0.37$ (mean $+0.24$); all 24 clear $|z| \ge 3.54$ analytically and
under the generation bootstrap, 21 of 24 have a cell-level $|z| \ge 2$, and 4 of 24 clear 3.54
under the cell-level bootstrap with its 15 or so clusters. Both bootstraps use a fixed seed (20261111).

\paragraph{Stratified standard errors and every control cell.} The analytic standard error of the marker mean, $s/\sqrt{n}$ over all of a cell's marker generations, mixes the systematic differences between lengths, spacings, and alphabets into the sampling variance. Treating those (length, spacing, alphabet) cells as fixed strata instead gives $\mathrm{SE}^2(\overline{\mathrm{acc}}_{\text{Bounds}}) = \sum_c w_c^2\, s_c^2/n_c$, with $w_c$ the share of marker generations in stratum $c$, and the same over alphabets within each length for the baseline. The stratified $z$ is larger than the analytic one in 44 of the 48 control cells, since the between-stratum spread no longer counts as noise, and gives the same verdicts at the four stronger models (23 of 24). Table~\ref{tab:controlcells} lists all 48 spaced-letter and digit cells with the length-matched baseline, the marker accuracy, $\Delta$, the share of remaining errors removed, and the three $z$ values.

\begin{table}[!htbp]
\centering
\scriptsize
\caption[Every pooled cell of the two tokenization controls ($N \le 64$, $k \in \{1, 2, 4, 8, 16\}$, all alphabet sizes).]{Every pooled cell of the two tokenization controls ($N \le 64$, $k \in \{1, 2, 4, 8, 16\}$, all alphabet sizes). No markers: the no-marker accuracy length-matched to the marker runs; share removed: $\Delta / (1 - \text{no-marker})$; $z$: generation-level standard errors (this report's default); $z$ (stratified): the fixed-strata estimator above; $z$ (cell bootstrap): resampling whole (length, spacing, alphabet) cells (open models only). Significance is $|z| \ge 3.54$.}
\label{tab:controlcells}
\adjustbox{max width=\linewidth,max totalheight=0.86\textheight}{\input{tables/control_cells.tex}}
\end{table}

We therefore read a pooled $z$ as a statement about the mean effect over the grid we ran, not as
evidence that the effect is uniform across lengths and spacings, nor that it would survive a
different choice of grid. The claims about structure rest on the length-resolved and
position-resolved analyses (Figures~\ref{fig:profile} and~\ref{fig:position}), not on the
pooled numbers.

\paragraph{The baseline's marker sentence.} The no-marker prompt keeps the sentence that
describes the markers. To measure what that costs, we reran the no-marker condition at the four
stronger models with every mention of markers removed from the prompt, on the same inputs: the ten task variants of the two tokenization
controls and the dependency family, $N \in \{16, 32, 64\}$, every alphabet or modulus value, 20
trials. Table~\ref{tab:cleanbase} gives the paired difference (clean minus with-sentence) per
model and family. It is $+0.006$ overall (95\% interval $[0.002, 0.009]$ over 300 cells), between
$-0.015$ and $+0.026$ in every model--family combination, and negative, meaning the sentence
helped, in 105 of the 300 cells. Recomputing Finding~2 at $N = 64$ against the clean baseline
moves the share of remaining errors removed by at most 0.014 at any model. Rerun at Qwen2.5-7B on the same cells, the sentence costs $-0.003 \pm 0.009$ (75 cells), so it does not confuse the smaller model either; over all 375 cells the cost is $+0.004$.

\begin{table}[!htbp]
\centering
\small
\caption[Cost of the baseline's marker sentence.]{Cost of the baseline's marker sentence: no-marker accuracy with the sentence removed minus with it present, on the same inputs, mean over (task, length, alphabet) cells with a 95\% interval over cells. $N \in \{16, 32, 64\}$, 20 trials per cell. The Qwen2.5-7B row was added for the review; the All row pools the five models (375 cells), against $+0.006$ over the 300 cells of the four stronger models quoted in the text.}
\label{tab:cleanbase}
\adjustbox{max width=\linewidth,max totalheight=0.86\textheight}{\input{tables/clean_baseline.tex}}
\end{table}

\paragraph{Dropping the redundant configurations.} At $A = 2$ the three position-local rules
coincide and at $M = 2$ \task{running\_sum\_local} computes \task{running\_xor}
(Chapter~\ref{sec:setup}). Recomputing all 91 pooled cells without those configurations moves
the effect by 0.013 in the median cell and at most 0.054; the effect gets \emph{larger} in 54 of
91 cells, because the two-symbol cells are the easiest and dilute the average. Three cells change
sign, all with $|\Delta| \le 0.009$ and none significant either way, and the Bonferroni verdict
changes in 5 of 91 cells. Finding~2 is unchanged: 23 of its 24 cells remain significant.

\paragraph{Paired inputs.} The main grids draw a fresh input for every (length, spacing, trial) cell, so the no-marker and marker runs of a trial transform different strings and every effect is an unpaired difference of means. To measure what that costs, we reran the spaced-letter and digit controls and the four dependency tasks at $N \in \{16, 32, 64\}$ with the input of each trial fixed across spacings (20 trials; Table~\ref{tab:paired}). The paired and unpaired estimates of the effect agree to within 0.01 in every model and format, the correlation between a trial's marker and no-marker accuracies is 0.16 to 0.35, and the paired standard error of a cell is 78\% to 91\% of the unpaired one, so pairing would have tightened the intervals by 9\% to 22\%. The paired run's no-marker rows reuse the main grid's seeds and reproduce its outputs exactly in 84\% to 100\% of generations, depending on the model.

\begin{table}[!htbp]
\centering
\small
\caption[Paired-input rerun.]{Paired-input rerun. Per model and format, averaged over cells (task, alphabet or modulus, $N$, $k > 0$): the marker effect from the main grid (different inputs per condition) and from the paired run (same inputs), the mean standard error of a cell under the unpaired and the paired estimator, and the mean correlation $r$ between a trial's marker and no-marker accuracies.}
\label{tab:paired}
\adjustbox{max width=\linewidth,max totalheight=0.86\textheight}{\input{tables/paired.tex}}
\end{table}

\section{Per-model results}
\label{app:models}

\paragraph{Mistral-7B.} In all three position-local tasks, the benefit is largest at short lengths and decays to zero at the longest length tested (\task{cyclic\_shift}, whose Mistral grid was extended to $N = 256$ with 40 trials per cell: at $N = 256$, $+0.008$, $z = 0.4$; \task{random\_permutation}, $N = 64$: $-0.009$, $z = 0.5$; \task{conditional\_shift}, $N=64$: $-0.006$, $z = 0.4$). In Qwen models the benefit holds or grows with length.

\paragraph{Repetition collapse.} On \task{running\_xor}, Qwen2.5-7B repeats a single wrong digit in about 38\% of trials, more often under dense markers. This failure falls to 0.9\% at 14B and 0.1\% for Mistral-7B, so it is specific to Qwen2.5-7B.

\paragraph{Out-of-alphabet answers.} On the digit tasks, Qwen2.5-7B often writes a letter (usually \texttt{t}, sometimes \texttt{o}) in place of \texttt{0}, in 34\% to 39\% of outputs, with or without markers. On spaced letters, several models step outside the alphabet (for example \texttt{I} after \texttt{H} instead of wrapping to \texttt{A}). Such characters are scored as errors.

\section{Per-position accuracy}
\label{app:perpos}

\paragraph{Pooled effects on dependency tasks.} The 8 significantly positive dependency cells of Chapter~\ref{sec:blindspot} are all on \task{running\_xor}, \task{running\_sum\_local}, or \task{sliding\_window\_max} at the stronger open models. Averaged over tasks, the effect ranges from $-0.05$ (Gemini 3.5) and $-0.04$ (Gemini 3.6 and Qwen2.5-7B) to $+0.04$ (Qwen2.5-32B). \task{copy\_previous}, which needs no computation, is negative even where the model is near ceiling without markers (Qwen3.8: $-0.26$). The pattern with length is consistent with this mixture: short sequences admit only dense spacings ($k < N$), so a larger share of their positions sit right after a marker, while long sequences add sparse spacings and are where the within-segment gain is large. At Qwen2.5-32B, \task{running\_xor} goes from $-0.11$ at $N = 4$ to $+0.13$ at $N = 64$, and \task{running\_sum\_local} from $-0.05$ to $+0.14$; Gemma 4's \task{running\_sum\_local} goes from $-0.10$ at $N = 8$ to $+0.19$ at $N = 64$. Qwen3.8 gains most on balance on \task{running\_xor} and \task{running\_sum\_local} ($+0.10$ and $+0.13$): its post-marker positions fail as at every model, but its within-segment positions on those two tasks gain $+0.13$ and $+0.14$.

Table~\ref{tab:perpos} gives the per-position means plotted in Figure~\ref{fig:position}. Table~\ref{tab:aligned} repeats the dependency-task comparison under alignment-robust scoring, Table~\ref{tab:blindk} repeats it at $k = 2$, 3, and 4 separately. ``No markers'' is the no-marker run's accuracy at the same position indices. Means are over the tasks in each block; Table~\ref{tab:blindwindow} splits \task{sliding\_window\_max} into the first and second positions after a marker and the strictly internal ones.

\begin{table}[!htbp]
\centering
\small
\caption[Per-position accuracy by model ($N \le 64$, $k \in \{2,4,8,16\}$), mean over the tasks in each block.]{Per-position accuracy by model ($N \le 64$, $k \in \{2,4,8,16\}$), mean over the tasks in each block.}
\label{tab:perpos}
\resizebox{0.85\linewidth}{!}{\input{tables/perpos.tex}}
\end{table}

\begin{table}[!htbp]
\centering
\small
\caption[\task{sliding\_window\_max} reads the two previous inputs, so the first two positions after a marker cross it.]{\task{sliding\_window\_max} reads the two previous inputs, so the first two positions after a marker cross it. Change against the no-marker run at the same positions for the first position after a marker ($i \bmod k = 0$), the second ($i \bmod k = 1$), and the strictly internal positions ($i \bmod k \ge 2$, only at $k \ge 3$; ---: none at $k = 2$), by spacing and pooled over $k \in \{2, 4, 8, 16\}$. Positions 0 and 1 of the sequence, which follow no marker, are in none of the classes.}
\label{tab:blindwindow}
\resizebox{0.85\linewidth}{!}{\input{tables/blind_spot_window.tex}}
\end{table}

\begin{table}[!htbp]
\centering
\small
\caption[Per-position gains on dependency tasks (marker run minus the no-marker run at the same positions) under index-by-index scoring and under edit-distance alignment, which repairs dropped or inserted characters.]{Per-position gains on dependency tasks (marker run minus the no-marker run at the same positions) under index-by-index scoring and under edit-distance alignment, which repairs dropped or inserted characters. $N \le 64$, $k \in \{2,4,8,16\}$. Share: the alignment-sensitive part of the within-segment gain, (index $-$ aligned) / index, shown only where the index-scored gain is positive; where that gain is tiny (the 7B-class models) the ratio is unstable and is not interpreted.}
\label{tab:aligned}
\resizebox{0.85\linewidth}{!}{\input{tables/aligned.tex}}
\end{table}

\paragraph{Error clustering.} Table~\ref{tab:autocorr} asks whether markers change how errors cluster along the sequence. For every generation at $N = 64$ of the four stronger open models, the per-position error indicator is computed under index scoring and under the edit-distance alignment, and three quantities are pooled over generations: the error rate, the probability $p_{1|1}$ that the position after an error is also wrong, and the probability $p_{1|0}$ that the position after a correct one is wrong. Under index scoring a dropped or inserted character makes every later position wrong, so $p_{1|1}$ there measures cascading misalignment; under alignment it measures the clustering of genuine local errors. Without markers the index-scored $p_{1|1}$ is 0.70 to 0.77 on the spaced and digit formats; sparse markers ($k \in \{2,4,8,16\}$) bring it to 0.24 to 0.59, and the mean run of consecutive errors shortens from 3.1 to 3.9 positions to 1.3 to 2.3. Under alignment the change is smaller (0.34 to 0.51 against 0.21 to 0.40) and $p_{1|0}$ falls together with the error rate, so sparse markers mostly stop a slip from propagating, with a smaller reduction in the clustering of substitutions. A marker after every character lowers the error rate but does not shorten the cascades (index-scored $p_{1|1}$ of 0.45 to 0.70). On the dependency tasks the same shift appears (0.63 to 0.79 against 0.40 to 0.61, index-scored), on top of the errors the blind spot adds at the post-marker positions. Read as a two-state Markov error model, $p_{1|1}$ and $p_{1|0}$ are its cascade and fresh-error parameters, and markers lower the first far more than the second.

\begin{table}[!htbp]
\centering
\small
\caption[Error clustering at $N = 64$ at the four stronger open models, pooled over the tasks of each format and over alphabets or moduli.]{Error clustering at $N = 64$ at the four stronger open models, pooled over the tasks of each format and over alphabets or moduli: error rate, $p_{1|1}$ (an error follows an error), and $p_{1|0}$ (an error follows a correct position), by marker density, under index scoring (top) and edit-distance alignment (bottom), and the mean length of a run of consecutive errors. Markers and spaces are stripped before scoring, so pairs of adjacent positions that straddle a marker are included; a run that reaches the end of the answer is counted at its length within the answer, without a correction for censoring, and under index scoring a short answer ends in one run to the end.}
\label{tab:autocorr}
\resizebox{0.85\linewidth}{!}{\input{tables/autocorr.tex}}
\end{table}

\begin{table}[!htbp]
\centering
\small
\caption[The blind spot at $k = 2$ to 6 separately (dependency tasks, $N \le 64$.]{The blind spot at $k = 2$ to 6 separately (dependency tasks, $N \le 64$; $k = 5$ and 6 were run at Qwen2.5-7B and 32B only): change at the post-marker and within-segment positions against the no-marker run at the same positions. $k = 3$ is otherwise excluded from the pooled analyses. At $k = 2$, every within-segment position of \task{sliding\_window\_max} is the second position after a marker, whose window still crosses it (Table~\ref{tab:blindwindow}).}
\label{tab:blindk}
\resizebox{0.85\linewidth}{!}{\input{tables/blind_spot_by_k.tex}}
\end{table}

\section{Reset variants}
\label{app:reset}

Figure~\ref{fig:reset} gives the per-position comparison of Section~\ref{sec:reset}. Pooled reset effects are not reported, because the reset prompt alone shifts no-marker accuracy by $-0.08$ to $+0.06$, so they would compare different baselines.

\begin{figure}[!htb]
\centering
\includegraphics[width=\linewidth]{figures/fig_reset_per_position.pdf}
\caption[Reset variants, per position.]{Reset variants, per position. Top row: accuracy at the position right after a marker, with the dependency intact (hollow circle) and with the recurrence restarting at each marker (arrow head, with its 95\% interval). Bottom row: the same for within-segment positions. Models: Qwen2.5-7B and 32B, Qwen3.8-27B, Gemma-4-31B; \task{sliding\_window\_max} was not run in reset form at 7B. $N \le 64$, $k \in \{2,4,8,16\}$; 95\% intervals.}
\label{fig:reset}
\end{figure}

\FloatBarrier
\section{External tasks, naturalistic tasks, and Gemini}
\label{app:external}

\paragraph{Naturalistic prefix-state tasks.} Two further tasks we built with natural readings, the nesting depth of a parenthesis string (modulo 4) and a running checksum of a digit string (modulo 10), carry the whole prefix as state, like the CLRS running minimum, rather than one previous input (Appendix Table~\ref{tab:naturalistic}). Qwen2.5-7B has low baseline accuracy on both (0.30 and 0.22 without markers, against chance levels of 0.25 and 0.10), so its cells are not interpretable. At Qwen2.5-32B and Gemma 4 markers hurt both (pooled $-0.00$ to $-0.14$; at Qwen2.5-32B on the checksum the pooled effect is near zero only because $k = 1$, where every position follows a marker, gains $+0.26$ while every $k \ge 2$ loses 0.02 to 0.15): the post-marker position falls by 0.07 to 0.21, and because the state carries forward, the within-segment positions fall with it (by 0.06 to 0.20) instead of gaining as they do on the one-step tasks. In three of the four task--model cells the digit written after a marker matches the answer of a state restarted at the marker more often than the true answer (0.41 to 0.51 against 0.23 to 0.37); only Gemma 4 on the checksum matches neither (0.16 against 0.32). The blind spot is confined to one position only when one position is all that needs what lies before the marker.

Tables~\ref{tab:external} and~\ref{tab:gemini} give the results on tasks we did not design (Chapter~\ref{sec:blindspot}) and on Gemini.

\begin{table}[!htbp]
\centering
\small
\caption[External position-local tasks, pooled across length and density.]{External position-local tasks, pooled across length and density. Base: no-marker accuracy averaged over lengths. Best-$k$: accuracy at the single best spacing (optimistic). $\Delta$: marker runs minus the length-matched no-marker mean (Chapter~\ref{sec:setup}). ProcBench-Substitute was regenerated after we made the vendored ProcBench generator's substitution table deterministic (it was drawn in a process-dependent order).}
\label{tab:external}
\adjustbox{max width=\linewidth,max totalheight=0.86\textheight}{\input{tables/external.tex}}
\end{table}

\begin{table}[!htbp]
\centering
\small
\caption[Closed-model check and CLRS-Text.]{Closed-model check and CLRS-Text. Gemini 3.5 Flash (3 trials per cell, letter and dependency tasks, pooled over all tested lengths, including $N$ up to 256 for \task{cyclic\_shift}) and Gemini 3.6 Flash (5 trials per cell, all 13 task--format variants, $N \le 64$), with the no-marker mean length-matched. CLRS-Text tasks (Gemini 3.5) use token-level markers and scoring.}
\label{tab:gemini}
\adjustbox{max width=\linewidth,max totalheight=0.86\textheight}{\input{tables/gemini.tex}}
\end{table}

We also ran CLRS-Text's Kadane algorithm; its base accuracy at both Qwen scales is statistically indistinguishable from chance under a shuffle test, so it is not reported in the tables or figures. Minimum is well above chance at Qwen2.5-32B ($z = 4.1$ from shuffled chance).

\paragraph{Gemini's tokenization, from the API's token counts.} We cannot inspect Gemini's tokenizer, but every response reports the prompt's token count, and across the Gemini 3.6 grid the instruction is fixed while the input grows. Table~\ref{tab:geminitokens} gives the prompt tokens added per added input character from $N = 32$ to $N = 64$, counting each marker as one token; the steps $4 \to 8$, $8 \to 16$, and $16 \to 32$ agree to within $0.04$. In the spaced and digit formats and on the binary dependency tasks, every added character adds exactly one token and every marker exactly one more, at every length and spacing, so the one-token-per-character property holds at Gemini 3.6 as at the open models. In the letter format, letters without markers cost $0.46$ tokens each (about 2.2 letters per token); a marker after every letter raises this to one token per letter, while markers every two, four, or eight letters leave it at $0.46$ to $0.51$, so at Gemini a sparse marker splits few of its letter tokens, which fits spacing beating markers there (Chapter~\ref{sec:tokenization}). The counts fix how many tokens the characters occupy, not which characters share one; the open-model verification by inspection remains the primary one.

\begin{table}[!htbp]
\centering
\small
\caption[Gemini 3.6 Flash.]{Gemini 3.6 Flash: prompt tokens added per added input character, from the API's token counts, by input format and marker spacing $k$, at the step from $N = 32$ to $N = 64$, with each marker counted as one token (exact in the spaced, digit, and binary formats, where the count at spacing $k$ exceeds the no-marker count by the number of markers at every $N$). A value of 1.00 is one token per character.}
\label{tab:geminitokens}
\adjustbox{max width=\linewidth,max totalheight=0.86\textheight}{\input{tables/gemini_tokens.tex}}
\end{table}

\paragraph{Naturalistic prefix-state tasks.} Two tasks built for this check have natural readings and carry the whole prefix as state: \task{bracket\_depth} (the input is a random string of parentheses; the output at each position is the nesting depth so far, modulo 4) and \task{running\_checksum} (the input is a digit string; the output at each position is the running sum modulo 10). They were run with the standard grid ($N \le 64$, all spacings, 20 trials per cell) at Qwen2.5-7B, Qwen2.5-32B, and Gemma 4. Table~\ref{tab:naturalistic} gives the pooled effect and the per-position comparison, with the four synthetic dependency tasks summarized the same way for the same models, and in the last four columns, over the marked generations ($k \ge 2$), the share of characters that match the true answer and the share that match the answer a model would give if its running state restarted at every marker; the two coincide whenever the state before the marker is 0 modulo the task's modulus, and such positions count for both.

\begin{table}[!htbp]
\centering
\small
\caption[Naturalistic prefix-state tasks.]{Naturalistic prefix-state tasks. No-marker accuracy, the pooled length-matched marker effect, the change at post-marker and within-segment positions against the no-marker run at the same positions ($k \in \{2, 4, 8, 16\}$, index scoring), and, over the marked generations with $k \ge 2$, the share of post-marker and within-segment characters that match the true answer and the share that match the answer of a state restarted at every marker. Positions where the two answers coincide count for both.}
\label{tab:naturalistic}
\adjustbox{max width=\linewidth,max totalheight=0.86\textheight}{\input{tables/naturalistic.tex}}
\end{table}

\section{Real sequences, all eight rules}
\label{app:real}

Table~\ref{tab:realfull} gives all eight real-sequence rules of Chapter~\ref{sec:real}. In the contiguous format the same-class, later-letter, and counting rules are near or below chance without markers at every model (chance is 0.5 for the two binary rules and 0.25 for the two counts modulo 4), so their marker effects there are not interpretable; the copy-previous rules were added after this was seen. With spaced characters, Gemma 4 is above chance on all four and Qwen2.5-32B on the same-class rule; Chapter~\ref{sec:real} discusses those cells.

\begin{table}[!htb]
\centering
\small
\caption[All eight real-sequence rules.]{All eight real-sequence rules; columns as in Table~\ref{tab:real}.}
\label{tab:realfull}
\adjustbox{max width=\linewidth,max totalheight=0.86\textheight}{\input{tables/real_data_full.tex}}
\end{table}


\section{Mechanism details}
\label{app:mechanism}

\paragraph{Head-count dose response.} To test whether the benefit is spread over heads, we repeated the ablation at $N = 64$ and $k = 8$ with the top 50 and top 200 marker-attending heads (2\% and 8\% of the model's heads), each with random-head and sham controls of the same size (Figure~\ref{fig:mechanism}c). Blocking 50 heads removes 1\% of the gain on spaced letters and 4\% on digits. Blocking 200 removes 28\% (paired drop $+0.155 \pm 0.072$) and 21\% ($+0.115 \pm 0.083$), while 200 random heads blocked from the markers, or the top 200 blocked from as many random non-marker tokens, remove nothing (drops of $-0.007$ to $-0.002$). Because both controls of that size are null, the drop is specific to removing attention to markers rather than a general disruption. The benefit thus runs partly through attention to the markers, but thinly spread: even 200 heads carry about a quarter of it. (On spaced letters, the drops of about 0.04 for the 10-head conditions and the 50-head controls all come from one trial, the same one in each case; without it they are within 0.01 of zero.)

\paragraph{Activation patching.}
\label{app:patching}
To ask where the blind spot is represented, we teacher-force Qwen2.5-32B on \task{running\_xor} and \task{running\_sum\_local} ($M = 4$) at $N = 32$, $k = 4$, 20 inputs per task, twice: with markers (the marked run) and without (the clean run), each prompt followed by its correct answer. At every post-marker answer position we read the next-token distribution and record the log-odds of the correct digit against the reset digit, the input digit at that position, which is the answer a model gives when it treats the segment as a fresh start, and whether the correct digit is the argmax. Qwen's tokenizer gives one token per digit and per bar, so answer positions align between the two runs once the bars are skipped. We cache the clean run's residual stream after each decoder layer and, in the marked run, overwrite the residual stream at one layer, or at every layer at once, at a chosen set of positions with the clean run's values at the corresponding positions: the input digit before each marker, the input digit after it, every input digit, or the answer digit written just before the marker. The log-odds are averaged over the post-marker positions where the correct and reset digits differ, which on both tasks excludes the positions whose previous input is 0, where the two answers coincide: 69 of the 140 post-marker positions remain on \task{running\_xor} (1 to 5 per input) and 100 of 140 on \task{running\_sum\_local} (3 to 7 per input); accuracy counts every post-marker position. Each input contributes the mean over its usable positions, the 20 inputs are averaged, and recovery is computed on those means: the patched log-odds minus the marked ones, divided by the clean minus the marked, so 0 is the marked run and 1 the clean run. A ratio is only meaningful when the clean--marked gap is large; it is 7 to 9 nats at Qwen2.5-32B. Table~\ref{tab:patching} gives the results, and Figure~\ref{fig:mechlayers} shows the recovery by layer at Qwen2.5-32B. Without patching, the marked run prefers the reset digit by 6 to 7 nats at the post-marker step and is at or near chance there (0.46 on \task{running\_xor}, 0.32 on \task{running\_sum\_local}, against 0.79 and 0.47 in the clean run). Patching every input digit at every layer recovers 75\% and 86\% of the log-odds gap; the input digit right after the marker alone recovers 59\% and 64\%; the digit before the marker recovers 12\% and 1\%, and the answer digit already written 18\% and 20\%. Patched one layer at a time, the recovery from the input digits rises through the first half of the network, peaks at 81\% and 94\% around layer 34 of 64, and falls to zero by layer 50; the post-marker input digit alone peaks at 33\% and 63\% at the same depth. Accuracy recovers less (up to 29\%), because the patched log-odds end near zero rather than clearly positive. The most influential site for the failure is therefore the residual-stream representation of the input tokens, chiefly of the first input of the segment, with the strongest single-layer recovery around layer 34 and none after layer 50; the answer already written carries little of it. Because the clean and marked runs differ in token positions and context, patching identifies where the difference acts on the prediction, not where it originates. At Qwen2.5-7B the same procedure is not informative: its teacher-forced clean run is itself at chance at $N = 32$ (0.49 on \task{running\_xor}, 0.28 on \task{running\_sum\_local}), the clean--marked gap is 0.4 to 1.1 nats rather than 7 to 9, and the recovery ratios swing between $-0.5$ and $1.1$ (Table~\ref{tab:patching}); we report them for completeness and do not interpret them, and the figure omits that model. At Gemma 4 (gaps of 10 and 13 nats) and Qwen3.8 (2.2 nats) the ordering of sites is the one found at Qwen2.5-32B: every input digit 60\% and 61\% and 46\% and 74\%, the input digit after the marker alone 48\% and 30\% at Gemma 4 and 48\% and 69\% at Qwen3.8, the digit before the marker at most 6\% and 21\%, the answer already written 15\% to 20\% and 8\% to 35\%, with single-layer peaks of the every-input curves around layer 26 of 60 at Gemma 4 and layers 25 to 34 at Qwen3.8 and nothing after layer 50. The best single layer for the input after the marker at Gemma 4 on \task{running\_xor} is layer 0 (0.33): patching that input right after the first layer recovers a third of the gap, recovery then dips, and a second band of 0.25 to 0.29 follows around layers 24 to 27; it is an early-layer value, not an indexing artifact. Where the single-layer recovery from the input digit after the marker falls from one layer to the next, the layer at which it falls is where later positions read that digit, so the two hybrid models show which kind of layer does the reading. At Qwen3.8 it is mostly the softmax-attention layers: its 16 softmax layers account for 62\% of the total decline over the two tasks (0.79 of 1.27 in recovery units), about five times as much per layer as its 48 linear-attention layers, and the largest single drop (0.19, \task{running\_sum\_local}) is at softmax layer 51. At Gemma 4, local and global layers contribute about equally per layer (0.020 and 0.018), as expected, because its 1,024-token local window covers our whole prompts, so the local--global split cannot matter at these lengths. With Qwen3.8's small clean--marked gap, this comparison is indicative only.

\begin{table}[!htbp]
\centering
\small
\caption[Activation patching at four models.]{Activation patching at four models: post-marker accuracy in the clean and marked runs, and the log-odds recovery when each site is patched at every layer, with the best single layer in parentheses.}
\label{tab:patching}
\adjustbox{max width=\linewidth,max totalheight=0.86\textheight}{\input{tables/patching.tex}}
\end{table}

\paragraph{Attention to the needed input.} A separate run reruns the teacher-forced inputs of the patching experiment at Qwen2.5-32B with attention weights kept. For the query position that predicts answer character $j$ (the token before it, which in the marked run is the output marker at a post-marker step) it records, per layer, the attention placed on input character $j-1$, on input character $j$, on the marker before input character $j$, and on answer character $j-1$, averaged over heads, and the largest single-head value on input $j-1$. Positions are post-marker ($j > 0$, $j \bmod k = 0$) or within-segment; the unmarked run is read at the same indices. Table~\ref{tab:attention} sums the head-averaged values over layers; Figure~\ref{fig:mechlayers} shows the strongest head by layer.

\begin{table}[!htbp]
\centering
\small
\caption[Attention from the step that writes answer character $j$ ($N = 32$, $k = 4$, 20 inputs per task), averaged over heads and summed over layers (64 at Qwen2.5-32B, 60 at Gemma 4, and the 16 softmax-attention layers of Qwen3.8, so sums compare conditions within a model, not across models), with the layer where the attention to input $j-1$ peaks.]{Attention from the step that writes answer character $j$ ($N = 32$, $k = 4$, 20 inputs per task), averaged over heads and summed over layers (64 at Qwen2.5-32B, 60 at Gemma 4, and the 16 softmax-attention layers of Qwen3.8, so sums compare conditions within a model, not across models), with the layer where the attention to input $j-1$ peaks. Clean: no markers; marked: markers in input and answer, so the step after a marker is read at the output bar; input markers only: marked input with an unmarked answer, so that step is read at the previous answer digit. The within-segment retention of 89\% to 96\% quoted in the text refers to the lookback band in layers 35 to 45 at Qwen2.5-32B, not to these sums.}
\label{tab:attention}
\adjustbox{max width=\linewidth,max totalheight=0.86\textheight}{\input{tables/attention_crossing.tex}}
\end{table}

\begin{figure}[!htbp]
\centering
\includegraphics[width=\linewidth]{figures/fig2_mechanism_merged.pdf}
\caption[Mechanism (Qwen2.5-32B).]{Mechanism (Qwen2.5-32B). (a) Total attention mass on all marker tokens rises with density at every length (\task{cyclic\_shift}, $A=2$); the number of markers rises faster, so this is not evidence that individual markers draw more attention than content tokens. (b) The most marker-attending head (layer 21, head 1) attends from each position to an adjacent marker. (c) Head-count dose response (\task{random\_permutation}, $A = 8$, $N = 64$, $k = 8$, 20 trials): share of the marker gain removed by blocking the top 10, 50, or 200 marker-attending heads from the markers, against as many random heads (random) and the top heads blocked from as many random non-marker tokens (sham), on spaced letters and digits; intervals are those of the drop divided by the gain. The wide spaced-letter intervals at 10 and 50 heads come from the single outlier trial described in Chapter~\ref{sec:mechanism}; at 200 heads the top heads remove 28\% and 21\% and the controls nothing. (d) \task{copy\_previous} at $N=32$: mean attention from each output position to the input token it needs (top-10 heads) is 0.535 at $k=1$, 0.101 at $k=2$, and 0.278 at $k=8$, matching accuracy (Base 0.83; $k=1$: 0.88; $k=2$: 0.65; $k=8$: 0.84).}
\label{fig:mechanism}
\end{figure}

\section{Marker ablations}
\label{app:marker}

\paragraph{Symbol, placement, and wording.} On contiguous-letter \task{cyclic\_shift} ($A = 2$, $k = 4$) at Qwen2.5-7B and 32B, where much of the effect is tokenization, the choice of symbol matters less than its familiarity. At $N = 64$ at Qwen2.5-32B, against 0.54 without markers, \texttt{|} reaches 0.76, \texttt{;} 0.66, and \texttt{\#} 0.69; a reserved special token that never appears in ordinary text (\texttt{<|object\_ref\_start|>}) falls far below the baseline (0.16), and a less frequent multi-token marker \texttt{<|>} placed every 12 characters, the same token budget as \texttt{|} every 4, reaches 0.71. Moving each marker by $\pm 1$ or $\pm 2$ positions, or placing the same number at random, falls below the baseline at most lengths at 7B but stays above it at 32B (0.73, 0.73, and 0.68 at $N = 64$). Formal and casual phrasings of the instruction help at 32B at every length ($+0.08$ to $+0.26$), while a terse cipher-style phrasing that says nothing about the markers beyond copying them hurts there ($-0.12$ to $+0.01$). These cover one task and two models; the dependency-task placement check below covers four tasks.

\paragraph{Jittered placement on the dependency tasks.} The placement ablation above concerns the letter-task gain. To ask whether the blind spot needs exact periodicity, we reran the four dependency tasks at $k = 4$ and 8 with every marker moved by $-1$, 0, or $+1$ positions from its periodic place (drawn independently per marker and trial), on the same inputs and seeds as the main grid's $k = 4$ and 8 cells, so each jittered generation is paired with a periodic one. Positions are classified from where the markers actually are. Table~\ref{tab:jitterdep} compares the two placements against the no-marker run at the same positions.

\begin{table}[!htbp]
\centering
\small
\caption[Periodic against jittered marker placement on the dependency tasks ($N \le 64$, $k \in \{4, 8\}$, same inputs.]{Periodic against jittered marker placement on the dependency tasks ($N \le 64$, $k \in \{4, 8\}$, same inputs; index scoring). Pooled $\Delta$ is against the length-matched no-marker mean; post-marker and within-segment changes are against the no-marker run at the same positions.}
\label{tab:jitterdep}
\adjustbox{max width=\linewidth,max totalheight=0.86\textheight}{\input{tables/jitter_dependency.tex}}
\end{table}

\FloatBarrier
\section{Delimiters, longer sequences, sampling, and larger state spaces}
\label{app:delimiter}

\paragraph{The constrained decoder of Section~\ref{sec:sides}.} Concretely, at each decoding step the processor decodes the answer so far, counts its non-marker characters $c$, and, when the condition calls for markers and $c$ is a positive multiple of $k$, masks every token except the marker; otherwise it masks every token containing a marker; once $c$ reaches $N$ it stops intervening. The masked set is every vocabulary token whose decoded string contains a bar, the forced token is the single-character bar, decoding is otherwise greedy with the grid's token budget, and generation ends at the model's end-of-turn token or the budget. The tasks, lengths, spacings, and trials are those of the instruction-based 2$\times$2. This changes what the model is allowed to generate, not only the format of the string it produces, so the same machinery is applied in all four cells. Decoding is greedy without beam search and the mask only removes tokens, so at every step the chosen token is the model's own highest-scoring allowed token; the mask cannot change the scores of the tokens it leaves.


\paragraph{Prompt delivery.} Qwen2.5's chat template inserts a system prompt of its own when none is given, so every run above carries one. Delivering the same prompts three other ways at Qwen2.5-32B leaves both effects in place (Appendix Table~\ref{tab:decoding}): with an empty system message, a task-oriented system message, or no chat template at all, the gain on \task{random\_permutation} with digits is $+0.36$ to $+0.39$ against $+0.37$ by default, and on \task{running\_xor} the post-marker position stays at chance (0.51 to 0.53, chance 0.5) while within-segment positions gain $+0.14$ to $+0.18$. These results suggest that the effects are robust to the prompt-delivery formats tested, including removing the chat template. The checkpoint remains instruction-tuned without its template; Appendix~\ref{app:base} compares a base checkpoint directly.

\paragraph{Larger state spaces.} The dependency tasks above use at most eight symbols. We reran \task{running\_sum\_local} at Qwen2.5-32B with modulus 8, 16, and 32, writing the symbols space-separated so that each is one token at every modulus. The post-marker position sits at chance at every modulus: 0.13 against a chance level of 0.125 at $M = 8$, 0.054 against 0.062 at $M = 16$, and 0.043 against 0.031 at $M = 32$, while within-segment positions gain $+0.03$ to $+0.08$ over the no-marker run at the same positions. The blind spot is therefore not an artifact of a small alphabet. One caveat bounds the largest case: without markers the model reaches only 0.075 at $M = 32$ against a chance level of 0.031, so that column is compressed toward the floor.

\paragraph{Sampled decoding.} All open-model runs above decode greedily (Gemini 3.6, which is sampled, shows the blind spot too; Chapter~\ref{sec:setup}). Repeating three tasks at Qwen2.5-32B with sampling ($N \in \{32, 64\}$, $k \in \{4, 16\}$; temperature 0.7 with top-$p$ 0.8, and temperature 1.0 with top-$p$ 0.95; one sample per prompt) changes little (Appendix Table~\ref{tab:decoding}). The post-marker position stays at chance, 0.51 to 0.52 on \task{running\_xor} and 0.26 to 0.27 on \task{running\_sum\_local} under all three decoding regimes, and the gain on \task{random\_permutation} with digits is $+0.37$ greedy, $+0.38$ at temperature 0.7, and $+0.42$ at temperature 1.0.

Table~\ref{tab:delimiter} gives the follow-up experiment with commas and line breaks as markers and sequences up to $N = 256$, and Table~\ref{tab:decoding} repeats three tasks with sampling instead of greedy decoding (Chapters~\ref{sec:segmentation} and~\ref{sec:blindspot}).

\begin{table}[!htbp]
\centering
\small
\caption[Delimiter type and longer sequences at Qwen2.5-32B ($k \in \{4, 16\}$ pooled, 20 trials per cell.]{Delimiter type and longer sequences at Qwen2.5-32B ($k \in \{4, 16\}$ pooled, 20 trials per cell; \task{random\_permutation} with digits at $A = 8$, \task{running\_sum\_local} at modulus 4). Change in accuracy: marker run minus no-marker run, pooled over positions. Post-marker and within-segment: marker run minus the no-marker run at the same positions. Each delimiter is named in its own constraint sentence. The \texttt{|} cells at $N = 64$ reproduce the main grid exactly (same prompts, greedy decoding). No output reached the generation budget.}
\label{tab:delimiter}
\adjustbox{max width=\linewidth,max totalheight=0.86\textheight}{\input{tables/delimiter.tex}}
\end{table}

\begin{table}[!htbp]
\centering
\small
\caption[Decoding and prompt delivery at Qwen2.5-32B ($N \in \{32, 64\}$, $k \in \{4, 16\}$ pooled, 20 prompts per cell.]{Decoding and prompt delivery at Qwen2.5-32B ($N \in \{32, 64\}$, $k \in \{4, 16\}$ pooled, 20 prompts per cell; \task{random\_permutation} at $A = 8$, \task{running\_sum\_local} at modulus 4). Main grid: greedy decoding with the chat template's own default system prompt, on the same prompts. Sampling: one sample per prompt. Empty system message: \texttt{""}, which suppresses the default; task-oriented: ``You are a precise, deterministic text transformer\ldots''; no chat template: the raw prompt as a completion. Change: marker run minus no-marker run, with a 95\% interval; post-marker and within-segment: against the no-marker run of the same condition at the same positions.}
\label{tab:decoding}
\adjustbox{max width=\linewidth,max totalheight=0.86\textheight}{\input{tables/decoding_delivery.tex}}
\end{table}

\FloatBarrier
\section{Base checkpoints, a finetuned separator, and a byte-level model}
\label{app:base}

\paragraph{Base checkpoint.} The main grids use instruction-tuned models through their chat templates, which a base checkpoint lacks. To compare the two on equal terms, we prompt the base and instruct versions of Qwen2.5-7B and Qwen2.5-32B the same way: no chat template, the task's instruction, three solved examples drawn from a seed stream disjoint from the grid (with markers at the same spacing as the query), and the query ($N \in \{16, 32, 64\}$, $k \in \{0, 2, 4, 8, 16\}$, 20 trials per cell; the four dependency tasks and the three digit tasks). Table~\ref{tab:base} summarizes by format. Both effects are present in the pretrained models. At 7B, on the digit tasks the base checkpoint gains $+0.10$ from markers (post-marker $+0.20$, within-segment $+0.06$) and the instruct version $+0.14$ ($+0.30$, $+0.11$); the few-shot prompt puts both 7B models well above their zero-shot grid accuracy on digits (0.55 and 0.48 without markers). On the dependency tasks both 7B models are near chance under this prompt (0.56 and 0.51), the post-marker change is $-0.07$ at the base checkpoint and $-0.03$ at the instruct one, and where the base checkpoint is clearly above chance, on \task{copy\_previous} (0.73), the position after a marker loses 0.20 while the rest of the segment loses 0.03. At 32B both checkpoints are at 0.96 on the digit tasks under the few-shot prompt and gain $+0.04$ each, so this pair says nothing about the long-sequence gain. On the dependency tasks, where both are far above chance (0.77 and 0.76 without markers), the blind spot is nearly as large in the pretrained model as after instruction tuning: the post-marker position loses 0.17 at the base checkpoint and 0.22 at the instruct one while the within-segment positions gain 0.02 and 0.03, and task by task the losses are 0.19 against 0.33 on \task{copy\_previous}, 0.31 against 0.33 on \task{running\_xor}, 0.16 against 0.17 on \task{running\_sum\_local}, and 0.03 against 0.05 on \task{sliding\_window\_max}. Instruction tuning enlarges both effects somewhat; it creates neither.

\begin{table}[!htbp]
\centering
\small
\caption[Base against instruction-tuned checkpoints of Qwen2.5-7B and Qwen2.5-32B under one prompt (no chat template, three solved examples.]{Base against instruction-tuned checkpoints of Qwen2.5-7B and Qwen2.5-32B under one prompt (no chat template, three solved examples; $N \in \{16, 32, 64\}$, $k \in \{2, 4, 8, 16\}$ against $k = 0$, 20 trials per cell). Digits: mean over the three digit tasks; Dependency: mean over the four dependency tasks. $\Delta$ is the length-matched pooled effect; per-position changes are against the no-marker run at the same positions.}
\label{tab:base}
\adjustbox{max width=\linewidth,max totalheight=0.86\textheight}{\input{tables/base_model.tex}}
\end{table}

\paragraph{A finetuned separator.} Can a small finetune teach the model to carry a dependency across a marker? We train a LoRA adapter (rank 16, all attention and MLP projections, two epochs, learning rate $2 \times 10^{-4}$, 16 sequences per step) on Qwen2.5-7B-Instruct with 2{,}040 examples of \task{running\_xor} and \task{running\_sum\_local} ($M \in \{2, 4, 8\}$) at $N \in \{8, 16, 32, 64\}$ and every spacing $k \in \{0, 2, 4, 8, 16\}$, with the marked answer as the target, so every crossing step is supervised (seeds disjoint from the grid). Two control adapters are trained on the same tasks with the 480 examples at $k = 0$ only, to separate learning the tasks from learning to cross a marker: one for two epochs, and one for nine, which matches the marked adapter's number of example passes (4{,}320 against 4{,}080). All three reach a training loss below $0.001$. The merged models are evaluated with the standard runners (Table~\ref{tab:lora}) on \task{running\_xor} at the grid's own inputs, which the adapter has not seen, on two dependency tasks that were not in training, and on a position-local digit task. On the trained tasks all three adapters are perfect without markers, but the controls' post-marker position falls: to 0.60 on \task{running\_xor} after two epochs ($-0.40$ against their own no-marker run) and to 0.56 after the matched nine ($-0.44$), and to 0.40 and 0.37 on \task{running\_sum\_local} ($-0.57$ and $-0.63$). Learning the task without markers widens the blind spot. The adapter trained on marked targets scores 1.00 at every position of both trained tasks. On \task{copy\_previous} it repairs the post-marker position ($+0.04$, from $-0.11$ in the released model; the matched control reaches $0.00$); on \task{sliding\_window\_max} it does not ($-0.07$, where the matched control reaches $+0.06$). The costs it shows on held-out tasks appear equally in the step-matched control: both lose no-marker accuracy on \task{sliding\_window\_max} (0.60 to 0.43 and 0.44) and both lose the marker gain on \task{random\_permutation} with digits ($+0.08$ in the released model; $-0.02$ and $-0.01$), so those costs come from the amount of finetuning rather than from marker supervision. What marker supervision does on its own is to remove the blind spot on the tasks it was shown, where the same training without markers widens it; on the held-out dependency tasks the step-matched control does as well or better at this size. Set against the worked-example prompt of Chapter~\ref{sec:blindspot}, which is a two-shot across-boundary demonstration, the finetune repairs the trained tasks completely where the prompt repaired them in part, and the prompt leaves other tasks untouched where the finetune does not. The two-epoch control saw 480 examples against 2{,}040, so the matched control is the fair comparison; the remaining asymmetry is that only the marked adapter ever saw a marker. Qwen2.5-7B is near chance on several of these tasks, so this is a bounded test of the idea.

\paragraph{The separator at Qwen2.5-32B.} The same recipe at Qwen2.5-32B-Instruct (one sequence per step with 16 accumulation steps, so the same 16 sequences per update; final training loss below $10^{-4}$ for every adapter) gives a marked adapter, a two-epoch unmarked control, and an unmarked control trained for nine epochs to match the marked adapter's number of example passes. On the trained tasks the result repeats: the marked adapter is perfect at every position, and both controls widen the blind spot ($-0.39$ to $-0.41$ on \task{running\_xor} and $-0.62$ to $-0.64$ on \task{running\_sum\_local}, against $-0.19$ and $-0.16$ released). On the held-out dependency tasks the step-matched control separates two cases. On \task{sliding\_window\_max} the repair is specific to marked training: the marked adapter reaches $+0.02$, against $-0.13$ released and $-0.27$ and $-0.23$ for the two controls. On \task{copy\_previous} it is not: the marked adapter reaches $+0.05$, but the step-matched control also reaches $0.00$ (two-epoch control $-0.22$). The cost in accuracy without markers likewise comes mostly from the number of training steps: the step-matched control lowers it on \task{copy\_previous} and on \task{random\_permutation} with digits (0.86 to 0.80 and 0.87 to 0.77), close to the marked adapter's 0.75 on both, while the shorter control raises it (0.95). On \task{random\_permutation} the marker gain survives in every adapter (post-marker $+0.25$ marked, $+0.32$ step-matched, $+0.15$ released).

\begin{table}[!htbp]
\centering
\small
\caption[LoRA separators at Qwen2.5-7B-Instruct and Qwen2.5-32B-Instruct, evaluated with the standard grid ($N \le 64$, 20 trials per cell).]{LoRA separators at Qwen2.5-7B-Instruct and Qwen2.5-32B-Instruct, evaluated with the standard grid ($N \le 64$, 20 trials per cell). ``Marked targets'': trained on \task{running\_xor} and \task{running\_sum\_local} at every spacing with marked answers; ``$k = 0$ only'': the same tasks without markers, for two epochs or for a number of steps matched to the marked adapter. Per task: no-marker accuracy (Base), and the change at the post-marker and within-segment positions against the checkpoint's own no-marker run at the same positions. The rows ``as released'' come from one 20-trial batch, so they differ slightly from the main grid's values, which pool 40 to 80 trials (for example $-0.16$ against $-0.15$ on \task{copy\_previous} at 32B).}
\label{tab:lora}
\adjustbox{max width=\linewidth,max totalheight=0.86\textheight}{\input{tables/lora.tex}}
\end{table}

\paragraph{A byte-level model.} EvaByte-SFT \citep{zheng2025evabyte}, a 6.5B byte-level chat model, gives one token per character and per marker; its released code needs one fix under current library versions (its rotary position tables are left uninitialized and must be rebuilt after loading). On \task{running\_xor}, \task{copy\_previous}, and two position-local digit tasks it writes the input back instead of applying the rule in 45\% of outputs and stays at or below chance with and without markers (every marker effect $|z| \le 1.8$), so it gives no test of either effect.

\section{Well-formedness and cost}
\label{app:wfr}

Character accuracy strips markers and spaces before comparing, so it does not say whether the
model followed the whole instruction. Table~\ref{tab:compliance} adds two stricter measures (with generation-weighted pooled rows) over
the 86{,}400 generations whose raw text we stored: \emph{exact transformation match}, the share whose output equals the answer
once markers and spaces are removed from both sides, and \emph{compliant}, the share whose raw first line carries markers
in exactly the right places (for no-marker runs, carries none). Weighted by generations, markers raise exact match on the position-local tasks, from 0.16 to 0.38 on contiguous letters (Qwen3.8 and Gemma 4 only), 0.39 to 0.43 on digits, and 0.39
to 0.42 on spaced letters, and lower it on the dependency tasks, from 0.24 to 0.11, the same
directions as character accuracy. Pooled over all task groups and weighted by the number of generations (the position-local controls contribute most), marker placement is followed in 95\% of marked generations at Gemma 4, 94\% at Qwen3.8, and 92\% at Qwen2.5-32B, against 60\% to 62\% at Mistral, Phi, and Qwen2.5-7B; without markers, models almost never emit one. On the dependency tasks alone compliance is lower, 0.94 at Gemma 4 and 0.89 at Qwen3.8 but 0.71 at Qwen2.5-32B and 0.65 at Qwen2.5-14B (at the Qwen2.5 models these rows cover \task{sliding\_window\_max} only, the dependency runs there that store raw output; Appendix~\ref{app:trials}), so at the Qwen2.5 models a substantial minority of marked outputs on those tasks do not carry the markers as asked. The blind spot is not a copying failure, however: in the constrained 2$\times$2 the output format is exact in the input-only and both-sides cells, and the post-marker position is at chance in both (Section~\ref{sec:sides}). The averages quoted in this paragraph are weighted by generations; the pooled rows of Table~\ref{tab:compliance} give them directly.

At Qwen2.5-32B (first seed batch), a well-formed output uses only the task alphabet and is at least as long as the answer; markers lower that share in six of seven tasks, by up to 0.11 (\task{conditional\_shift}, 0.86 to 0.75; \task{cyclic\_shift} 0.52 to 0.44, \task{running\_sum\_local} 0.77 to 0.68, \task{copy\_previous} 0.96 to 0.88). On \task{cyclic\_shift}, \task{copy\_previous}, and \task{sliding\_window\_max} the drop comes from the completion rate: more marked outputs stop short of $N$ answer characters (on \task{cyclic\_shift}, 0.92 complete without markers against 0.81 with), so markers that curb running past the required length (Table~\ref{tab:wallclock}) also make stopping too early somewhat more frequent. On \task{conditional\_shift} and \task{running\_sum\_local} markers also lower the share of outputs that use only the task alphabet.

\begin{table}[!htbp]
\centering
\small
\caption[Exact transformation match (markers and spaces removed from both sides) and marker-placement compliance by model ($N \le 64$, all spacings pooled).]{Exact transformation match (markers and spaces removed from both sides) and marker-placement compliance by model ($N \le 64$, all spacings pooled): the dependency tasks and all task groups weighted by generations, and every task group pooled over models. Exact: the cleaned output equals the answer. Compliant: the raw first line carries markers at exactly the expected offsets, or none for the no-marker runs. For Qwen2.5-7B, 14B, and 32B, Mistral, and Phi the ``Dependency, as is'' rows cover only \task{sliding\_window\_max} (at Phi also \task{copy\_previous}), the dependency runs of those models that store raw output, and the contiguous-letter rows cover only Qwen3.8 and Gemma 4, so the pooled letter row averages over those two models (Appendix~\ref{app:trials}).}
\label{tab:compliance}
\adjustbox{max width=\linewidth,max totalheight=0.86\textheight}{\input{tables/compliance_short.tex}}
\end{table}

\begin{table}[!htbp]
\centering
\small
\caption[Wall-clock seconds per call at Qwen2.5-32B on contiguous-letter \task{cyclic\_shift} ($A = 2$).]{Wall-clock seconds per call at Qwen2.5-32B on contiguous-letter \task{cyclic\_shift} ($A = 2$). Throughput is 14 to 15 tokens/s in every condition. No marker is inserted after the last character, so the $k = 8$ cell at $N = 8$ contains none and is omitted.}
\label{tab:wallclock}
\adjustbox{max width=\linewidth,max totalheight=0.86\textheight}{\begin{tabular}{lcccccc}
\toprule
Condition & $k$ & $N=8$ & 16 & 32 & 64 & 128 \\
\midrule
Base & -- & 0.330 & 0.570 & 1.270 & 3.763 & 11.380 \\
Bounds & 1 & 1.133 & 2.215 & 4.403 & 7.684 & 14.260 \\
Bounds & 2 & 0.544 & 1.076 & 2.142 & 4.213 & 9.000 \\
Bounds & 4 & 0.380 & 0.787 & 1.729 & 3.298 & 6.453 \\
Bounds & 8 & --- & 0.613 & 1.288 & 2.479 & 4.976 \\
\bottomrule
\end{tabular}}
\end{table}
\end{document}
